\documentclass[journal]{IEEEtran}

% Essential packages
\usepackage[utf8]{inputenc}
\usepackage[T1]{fontenc}
\usepackage{amsmath,amssymb,amsfonts}
\usepackage{graphicx}
\usepackage{booktabs}
\usepackage{tabularx}
\usepackage{array}
\usepackage{multirow}
\usepackage{siunitx}
\usepackage{microtype}
\usepackage{cite}
\usepackage{url}
\makeatletter
\@namedef{ver@fixltx2e.sty}{9999/12/31}
\makeatother
\usepackage{dblfloatfix}
\usepackage{etoolbox}

\AtBeginEnvironment{enumerate}{\setlength{\itemsep}{0pt}\setlength{\parsep}{0pt}\setlength{\topsep}{0.5pt plus 0.5pt minus 0.5pt}\setlength{\partopsep}{0pt}}
\AtBeginEnvironment{itemize}{\setlength{\itemsep}{0pt}\setlength{\parsep}{0pt}\setlength{\topsep}{0.5pt plus 0.5pt minus 0.5pt}\setlength{\partopsep}{0pt}}

% Modest, consistent display spacing without changing body leading or math size.
\newcommand{\manuscriptdisplayspacing}{%
  \setlength{\abovedisplayskip}{4.5pt plus 1pt minus 1pt}%
  \setlength{\belowdisplayskip}{4.5pt plus 1pt minus 1pt}%
  \setlength{\abovedisplayshortskip}{3.5pt plus 1pt minus 1pt}%
  \setlength{\belowdisplayshortskip}{3.5pt plus 1pt minus 1pt}%
  \setlength{\jot}{4pt}%
}
\AtBeginDocument{\manuscriptdisplayspacing}

% IEEEtran bibliography customization
\def\IEEEbibitemsep{0.5pt plus 0.2pt minus 0.1pt}

\usepackage[hidelinks]{hyperref}

% Custom column types for tabularx
\newcolumntype{Y}{>{\centering\arraybackslash}X}
\newcolumntype{L}[1]{>{\raggedright\arraybackslash}p{#1}}
\newcolumntype{C}[1]{>{\centering\arraybackslash}p{#1}}
\newcolumntype{R}[1]{>{\raggedleft\arraybackslash}p{#1}}

% Float placement parameters to keep page composition tight and compact
\renewcommand{\topfraction}{0.95}
\renewcommand{\bottomfraction}{0.95}
\renewcommand{\textfraction}{0.05}
\renewcommand{\floatpagefraction}{0.85}
\renewcommand{\dbltopfraction}{0.95}
\renewcommand{\dblfloatpagefraction}{0.85}
\setcounter{dbltopnumber}{2}
\setlength{\textfloatsep}{4.0pt plus 1pt minus 1pt}
\setlength{\floatsep}{3.5pt plus 1pt minus 1pt}
\setlength{\dbltextfloatsep}{4.0pt plus 1pt minus 1pt}
\setlength{\dblfloatsep}{3.5pt plus 1pt minus 1pt}
\setlength{\intextsep}{3.5pt plus 1pt minus 1pt}
\setlength{\abovecaptionskip}{2.0pt}
\setlength{\belowcaptionskip}{1.0pt}

% Allow line breaks in paths and URLs
\def\UrlBreaks{\do\/\do-\do_}

% Safe figure inclusion macro (checks local, ../figures/, and figures/ directories)
\graphicspath{{../figures/}{./figures/}{./}}
\newcommand{\safefig}[2][width=\linewidth]{%
  \IfFileExists{#2}{\includegraphics[#1]{#2}}{%
    \IfFileExists{../figures/#2}{\includegraphics[#1]{../figures/#2}}{%
      \IfFileExists{figures/#2}{\includegraphics[#1]{figures/#2}}{%
        \fbox{\begin{minipage}[c][4cm][c]{0.95\linewidth}\centering
        \textbf{Placeholder: Figure \texttt{#2} not found.}\\
        \textit{Please ensure \texttt{#2} is uploaded to the project directory.}
        \end{minipage}}%
      }%
    }%
  }%
}

% Redefine abstract and keywords environments so body text is rendered in normal font (not bold)
\makeatletter
\renewenvironment{abstract}{%
  \if@twocolumn
    \normalfont\@IEEEabskeysecsize\textbf{\textit{\abstractname}:}\hspace{0.5em}%
  \else
    \bgroup\par\addvspace{0.5\baselineskip}\centering\@IEEEabskeysecsize\textbf{\abstractname}\par\addvspace{0.5\baselineskip}\egroup\quotation\@IEEEabskeysecsize\mdseries
  \fi\@IEEEgobbleleadPARNLSP
}{\relax\ifCLASSOPTIONtechnote\vspace{1.34ex}\else\vspace{0.67ex}\fi
  \par\if@twocolumn\else\endquotation\fi\normalfont\normalsize}

\renewenvironment{IEEEkeywords}{%
  \if@twocolumn
    \normalfont\@IEEEabskeysecsize\textbf{\textit{\IEEEkeywordsname}:}\hspace{0.5em}%
  \else
    \bgroup\par\addvspace{0.5\baselineskip}\centering\@IEEEabskeysecsize\textbf{\IEEEkeywordsname}\par\addvspace{0.5\baselineskip}\egroup\quotation\@IEEEabskeysecsize\mdseries
  \fi\@IEEEgobbleleadPARNLSP
}{\relax\ifCLASSOPTIONtechnote\vspace{1.34ex}\else\vspace{0.67ex}\fi
  \par\if@twocolumn\else\endquotation\fi\normalfont\normalsize}
\makeatother

\title{Constrained Lifecycle Optimization of Low-Carbon Structural Concrete Under 11-Year Natural Carbonation Kinetics and Compressive Strength Demands}

\author{Abhitej~Singh~Bhullar, \IEEEmembership{Roll No.~23110009}\\
Department of Civil Engineering\\
\textit{Modeling of Earth System and Sustainability (MESS)}%
}

\markboth{Modeling of Earth System and Sustainability (MESS), 2026}{Bhullar: Constrained Lifecycle Optimization of Low-Carbon Structural Concrete}

\begin{document}

\maketitle

\begin{abstract}
Supplementary cementitious materials reduce emissions but deplete alkaline buffering against carbonation. Using an 11-year dataset of 45 mixtures, we formulate a constrained lifecycle optimization framework. With IS~10262 strength margins, discrete screening selects 50\% slag Mix~C24 at $\ge30\ \mathrm{mm}$ cover under EN~206; IS~456 Severe excludes C24, while C40/50 selects C23. At 25\,mm cover, kinetic uncertainty increases carbon by $16.5\%$. Across 18 scenarios, continuous optimization predicts $4.6\%\text{--}31.2\%$ carbon savings ($1.1\%\text{--}29.0\%$ with PCE). At 30\,mm cover, net global warming potential (GWP) is $133.4\ \mathrm{kg\ CO_2e/m^3}$, saving $14.1\ \mathrm{kg/m^3}$ over C24 ($8.8\text{--}9.9\ \mathrm{kg/m^3}$ with PCE). Uptake offsets 3.8\% to 5.7\% of emissions. No original mixture used PCE; chemical compaction to 2\% air remains uncalibrated.
\end{abstract}

\begin{IEEEkeywords}
Low-carbon concrete, natural carbonation kinetics, lifecycle optimization, supplementary cementitious materials, service-life durability, Pareto frontier.
\end{IEEEkeywords}

\section{Introduction}

\subsection{The Decarbonization Mandate for Structural Concrete}
Global warming potential (GWP) expresses greenhouse-gas impacts as carbon dioxide equivalent (CO$_2$e).

Annual concrete production exceeds 30 billion tonnes. Clinker calcination ($\mathrm{CaCO}_3 \to \mathrm{CaO} + \mathrm{CO}_2$) and fuel combustion contribute $7\%\text{--}8\%$ of global $\mathrm{CO}_2$ \cite{scrivener2018, monteiro2017}. Standard life-cycle inventories report benchmark OPC emissions of $0.832\text{--}0.840\ \mathrm{kg\ CO_2e/kg}$ \cite{ice3, scrivener2018}. Reported Indian kiln factors span $0.893\text{--}0.940\ \mathrm{kg\ CO_2e/kg}$ ($893\text{--}940\ \mathrm{kg/t}$) in plant audits and cradle-to-gate LCAs \cite{karadumpa2024}.

Supplementary cementitious materials (SCMs), including blast-furnace slag (slag), siliceous fly ash (FA), limestone filler (LS), and volcanic pozzolan (PO), replace clinker with industrial by-products \cite{scrivener2018, lothenbach2011}. Clinker replacement also creates a durability trade-off. Hydration of cement produces Portlandite ($\mathrm{Ca(OH)}_2$), which maintains pore solution $\mathrm{pH} > 12.5$ and passivates steel reinforcement through a microscopic oxide film ($\gamma\text{-}\mathrm{Fe}_2\mathrm{O}_3$). Atmospheric $\mathrm{CO}_2$ diffuses into unsaturated pores and neutralizes this alkaline reserve:
\begin{equation}
\mathrm{Ca(OH)}_2 + \mathrm{CO}_2 \to \mathrm{CaCO}_3 + \mathrm{H}_2\mathrm{O}
\label{eq:carbonation_rxn}
\end{equation}
This reaction lowers pore solution pH below 9.0 \cite{papadakis1991}. When the front reaches the reinforcement depth ($d_{\mathrm{cover}}$), steel depassivates, initiating active corrosion and cover spalling \cite{fib34, fib2020}. SCMs contain less reactive $\mathrm{CaO}$ and consume Portlandite pozzolanically. High clinker replacement therefore depletes the alkaline buffering capacity \cite{papadakis1992}.

\subsection{Research Gaps and Problem Formulation}
Accelerated carbonation tests ($1\%\text{--}4\%\ \mathrm{CO}_2$) alter pore microstructure relative to natural exposure \cite{vu2019, hule2024}. Prior optimizations treat strength or carbon in isolation, omitting reactive transport, entrapped air, and codified bounds.

This study uses the $\approx 11$-year Chennai dataset of 45 mixes \cite{hule2024, hule2024mendeley} to examine how clinker replacement can reduce lifecycle carbon while satisfying strength and carbonation-cover constraints. The framework separates field measurements, calibrated models, and literature-based scenarios into three evidence tiers. It combines a cross-validated, air-normalized Bolomey strength model with sheltered and unsheltered carbonation kinetics, including $+1\sigma$ margins, to screen measured mixtures and optimize continuous compositions across 18 scenarios. Closed-form durability bounds and shadow prices then explain how cover and composition limits control the selected mixtures. The notation is summarized in Table~\ref{tab:notation}.



\section{Decade-Long Natural Weathering Dataset and Transport Physics}

\subsection{Experimental Campaign and Field Exposure}
The experimental campaign comprises 45 concrete mixtures formulated with Portland cement, fly ash ($15\%\text{--}50\%$), slag ($30\%\text{--}50\%$), limestone ($10\%\text{--}20\%$), pozzolan ($15\%\text{--}30\%$), and ternary blends ($30\%\text{--}40\%$), exposed under sheltered and unsheltered conditions in Chennai, India \cite{hule2024, hule2024mendeley}. Water-to-binder ratios span $w/b \in [0.45, 0.60]$ with binder $B \in [220, 400]\ \mathrm{kg/m^3}$ (Mix C1 is the sole mix at $B = 220\ \mathrm{kg/m^3}$, $0\%$ SCM; SCM blends require $B \ge 250\ \mathrm{kg/m^3}$).

Concrete prisms were stored for 7--104 days (mean: 39.8 days; median: 28 days), with no statistically significant effect of storage duration observed ($p > 0.50$; Appendix~\ref{sec:appendix_standards}). The dataset contains 353 carbonation depth measurements across six testing periods spanning 1.0--11.5 years, including 167 sheltered and 186 unsheltered measurements. The final testing period had a mean exposure age of 10.75 years. Fresh concrete air content ranged from 0.9\% to 6.6\%, averaging 2.5\%. The study evaluates three target compressive strengths (30, 35, and 40 MPa) and three concrete cover depths (25, 30, and 35 mm) under both exposure conditions.

\begin{figure*}[!t]
\centering
\safefig[width=0.80\textwidth]{fig1_carbonation_kinetics_by_condition.png}
\caption{Longitudinal natural carbonation kinetics ($d$ vs. $\sqrt{t}$) across 11 years (final epoch 10.1--11.5\,yr, mean 10.75\,yr) of natural field exposure in Chennai, India: (a) Sheltered exposure (indoor/covered, EN 206 class XC3), showing accelerated penetration driven by dry convective $\mathrm{CO}_2$ diffusion; (b) Unsheltered outdoor exposure (rain-exposed, EN 206 class XC4), illustrating substantial carbonation attenuation ($26\%\text{--}30\%$ lower rates), consistent with physical pore-blocking from periodic monsoon wetting and pore condensation.}
\label{fig:kinetics_condition}
\end{figure*}

\begin{table}[!t]
\caption{Mathematical Notation, Parameters, and Physical Units}
\label{tab:notation}
\centering
\scriptsize
\renewcommand{\arraystretch}{0.55}
\setlength{\tabcolsep}{2.0pt}
\begin{tabularx}{\columnwidth}{l X c}
\toprule
\textbf{Symbol} & \textbf{Physical Meaning / Definition} & \textbf{Unit} \\
\midrule
$\mathbf{x}$ & Decision vector $[w/b, B, x_{\mathrm{FA}}, x_{\mathrm{SG}}, x_{\mathrm{LS}}, x_{\mathrm{PO}}]^T$ & --- \\
$w/b, B/W$ & Water-to-binder mass ratio and reciprocal $B/W$ & --- \\
$B$ & Total cementitious binder content & $\mathrm{kg/m^3}$ \\
$x_j, p_j$ & Mass fraction ($x_j$) and percentage ($p_j = 100 x_j$) of SCM $j$ & ---, \% \\
$\mathrm{Air}\%$ & Entrapped batch fresh air volume percentage & \% \\
$f_{c,28}, f_{\mathrm{target}}$ & 28-day cylinder compressive strength and structural demand & $\mathrm{MPa}$ \\
$\Delta f_{\mathrm{margin}}$ & Codified strength screening buffer ($+8.25\ \mathrm{MPa}$, IS~10262) & $\mathrm{MPa}$ \\
$x_c(t), k, d_{50}$ & Carbonation depth, rate ($x_c = k\sqrt{t}$), and 50-yr depth & $\mathrm{mm}, \mathrm{mm/yr^{0.5}}$ \\
$d_{\mathrm{cover}}, d_{50,\mathrm{allow}}$ & Reinforcement cover depth and allowable nominal depth & $\mathrm{mm}$ \\
$D_{\mathrm{eff}}, a, C_{\mathrm{ext}}$ & Effective $\mathrm{CO}_2$ diffusivity, buffer capacity, ambient $\mathrm{CO}_2$ & $\mathrm{m^2/s}, \mathrm{mol/m^3}$ \\
$S_w, \phi_p, \phi_g$ & Pore saturation ($0 \le S_w \le 1$), total and gas porosity & --- \\
$GWP_{\mathrm{net}}, GWP_{\mathrm{A1-A3}}$ & Net lifecycle GWP and cradle-to-gate embodied carbon & $\mathrm{kg\ CO_2e/m^3}$ \\
$\mathrm{CO}_{2,\mathrm{up}}^{\mathrm{B1}}$ & Use-phase 50-year stoichiometric carbonation uptake & $\mathrm{kg\ CO_2/m^3}$ \\
$EF_i, EF_{\mathrm{batching}}$ & Material emission factor ($i$) and batching footprint & $\mathrm{kg/kg}, \mathrm{kg/m^3}$ \\
$\mu_k$ & Karush-Kuhn-Tucker dual multiplier (shadow price) of constraint $k$ & $\mathrm{unit}^{-1}$ \\
\bottomrule
\end{tabularx}
\end{table}

\subsection{Earth System Reactive-Transport Governing Equations}
Concrete is an unsaturated geochemical porous medium undergoing reactive-diffusive exchange with atmospheric $\mathrm{CO}_2$. The 1D conservation law for gaseous $\mathrm{CO}_2$ concentration $C(x,t)$ [$\mathrm{mol/m^3}$] in unsaturated pores is:
\begin{equation}
\begin{aligned}
\frac{\partial (\phi_g C)}{\partial t}
&= \frac{\partial}{\partial x} \left( D_{\mathrm{eff}}(S_w, \phi_p, T) \frac{\partial C}{\partial x} \right) \\
&\quad - r_{\mathrm{carb}}
\end{aligned}
\label{eq:pde_transport}
\end{equation}
where $\phi_g = \phi_p (1 - S_w)$ is gas-filled porosity, $\phi_p$ is total connected porosity, $S_w \in [0,1]$ is liquid saturation, and $r_{\mathrm{carb}}$ is the volumetric mineralization sink ($\mathrm{mol\ CO_2 / (m^3\cdot s)}$).

Under the Millington-Quirk tortuosity-porosity formulation:
\begin{equation}
D_{\mathrm{eff}}(S_w) = D_0 \cdot \frac{\phi_g^{10/3}}{\phi_p^2} = D_0 \cdot \phi_p^{4/3} (1 - S_w)^{10/3}
\label{eq:millington_quirk}
\end{equation}
where $D_0 \approx 1.6 \times 10^{-5}\ \mathrm{m^2/s}$ is molecular $\mathrm{CO}_2$ diffusivity in air. In pore water, aqueous diffusivity decreases to $D_{\mathrm{aq}} \approx 1.8 \times 10^{-9}\ \mathrm{m^2/s} \approx 10^{-4} D_0$. In outdoor exposure (XC4), periodic monsoon rain induces capillary saturation ($S_w \to 1$), reducing $D_{\mathrm{eff}} \to 0$ and providing a plausible physical explanation for the $26\%\text{--}30\%$ empirical carbonation rate attenuation observed outdoors relative to sheltered exposure (XC3).

The assumptions are instantaneous reaction-front propagation ($r_{\mathrm{carb}} \to \infty$ at $x = x_c(t)$, M1), quasi-steady pore diffusion ($\partial^2 C / \partial x^2 \approx 0$ in $0 \le x \le x_c$, M2), and boundary values $C(0,t) = C_{\mathrm{ext}}$, $C(x_c, t) = 0$ (M3).

Integrating Laplace's equation under M3 gives a linear concentration gradient $\partial C / \partial x = -C_{\mathrm{ext}} / x_c(t)$. Applying the Stefan mass-conservation condition across the moving interface yields:
\begin{equation}
a \cdot \frac{\mathrm{d}x_c}{\mathrm{d}t} = -D_{\mathrm{eff}} \left. \frac{\partial C}{\partial x} \right|_{x = x_c^-} = D_{\mathrm{eff}} \frac{C_{\mathrm{ext}}}{x_c(t)}
\label{eq:stefan_front}
\end{equation}
where $a = [\mathrm{Ca(OH)}_2]_0 + \beta [\mathrm{C}\text{-}\mathrm{S}\text{-}\mathrm{H}]_0$ is total volumetric alkaline buffering capacity ($\mathrm{mol\ CO_2 / m^3}$). Separating variables and integrating from $x_c(0) = 0$:
\begin{equation}
\begin{aligned}
\int_0^{x_c(t)} x' \,\mathrm{d}x'
&= \frac{D_{\mathrm{eff}} C_{\mathrm{ext}}}{a} \int_0^t \mathrm{d}t' \\
\implies \frac{1}{2} x_c(t)^2
&= \frac{D_{\mathrm{eff}} C_{\mathrm{ext}}}{a} t
\end{aligned}
\label{eq:stefan_integral}
\end{equation}
Taking the square root recovers the exact Fickian rate law:
\begin{equation}
x_c(t) = k \sqrt{t}, \quad \text{with } k = \sqrt{\frac{2 D_{\mathrm{eff}} C_{\mathrm{ext}}}{a}}
\label{eq:fickian_k}
\end{equation}
Replacing clinker with SCMs reduces $a$, amplifying $k$. Slag hydration produces secondary C-S-H that retains alkaline reserve and refines pore structure, whereas fly ash pozzolanic reactions consume Portlandite while generating slow early microstructure, explaining why $\gamma_{\mathrm{FA}} / \gamma_{\mathrm{SG}} = 2.17\times$.

When continued hydration causes time-decaying diffusivity $D(t) = D_0 (t_0/t)^\alpha$ ($\alpha \in [0, 1)$), the moving boundary obeys anomalous sub-diffusive power-law penetration:
\begin{equation}
\begin{aligned}
x_c(t) &= \left[ \frac{2 D_0 C_{\mathrm{ext}} t_0^\alpha}{a (1 - \alpha)} \right]^{1/2} t^{(1 - \alpha)/2} \\
&= k' \, t^n, \quad n = \frac{1 - \alpha}{2}
\end{aligned}
\label{eq:power_law}
\end{equation}
In sheltered exposure, negligible post-curing hydration aging yields $\alpha \approx 0 \implies n \approx 0.50$ (empirical pooled sheltered fit $n = 0.509$, 95\% CI $[0.458, 0.560]$). Because $n = 0.509$ is slightly above the theoretical bound ($n = (1 - \alpha)/2 \le 0.50$), it is interpreted as $\alpha \approx 0$ within fitting error. For outdoor monsoon exposure, rain-induced late hydration yields $\alpha \approx 0.10 \implies n = 0.452 < 0.50$ (95\% CI $[0.406, 0.498]$).

\subsection{Authoritative Three-Tier Evidence Demarcation}
The analysis separates its inputs and outputs into three evidence tiers, summarized in Table~\ref{tab:evidence_hierarchy}. Tier 1 contains measured data, including carbonation depths, cylinder strengths, fresh air contents, and Chennai climate records. Tier 2 contains strength and kinetic models calibrated against the experimental cohort; the calibrated kinetic laws have $R^2 > 0.86$. These fitted relationships support interpolation within the observed mixture domain but do not make every optimized composition an experimentally tested mixture.

Tier 3 contains literature-based scenarios and extrapolations, including projected service-life depths and lifecycle carbon estimates. External inputs include EN~16757 uptake \cite{en16757} and EN~206 binder limits \cite{en206}. This distinction is retained when interpreting the optimization results: predicted performance at the design horizon and beyond the tested composition domain remains conditional on the adopted models and assumptions.

\begin{table*}[!t]
\caption{Authoritative Scientific Demarcation of Experimental and Model Variables Across Three Evidence Tiers}
\label{tab:evidence_hierarchy}
\centering
\scriptsize
\renewcommand{\arraystretch}{0.65}
\setlength{\tabcolsep}{3.0pt}
\begin{tabularx}{\textwidth}{L{2.2cm} c L{2.8cm} L{2.8cm} >{\raggedright\arraybackslash}X}
\toprule
\textbf{Variable / Quantity} & \textbf{Tier} & \textbf{Empirical Grounding} & \textbf{Model / External Basis} & \textbf{Defensible Scientific Claim} \\
\midrule
\textbf{Carbonation Depth ($x_c$)} & \textbf{Tier 1} & 353 phenolphthalein tests, 6 epochs (1.0--11.5\,yr, final epoch mean 10.75\,yr) & Indicator spraying; corner-filtered spatial average \cite{hule2024} & Direct field measurements over $\approx 11$ years in tropical Chennai. \\
\addlinespace[1pt]
\textbf{28-Day Strength ($f_{c,28}$)} & \textbf{Tier 1} & 45 cylinder tests per EN~12390-3 & Laboratory testing per standard protocol & Direct cylinder strength; characteristic targets $f_{ck,\mathrm{cyl}} = 30, 35, 40\ \mathrm{MPa}$ map to cube M40, M45, M50 per EN~206 / IS~456. \\
\addlinespace[1pt]
\textbf{Fresh Air Content ($\mathrm{Air}\%$)} & \textbf{Tier 1} & Batch air volume ($0.9\%\text{--}6.6\%$, mean $2.5\%$, median $1.8\%$) & Direct pressure gauge measurement & Entrapped air ($0.9\%\text{--}6.6\%$, Mix C24 was $1.6\%$) dictates compressive strength. \\
\addlinespace[1pt]
\textbf{Carbonation Rate ($k$)} & \textbf{Tier 2} & OLS origin regression across 45 mixes over $\approx 11$ years & $x_c = k\sqrt{t}$ (curing omitted as $p > 0.50$) & Mean sheltered $k = 3.73$, unsheltered $2.75\ \mathrm{mm/yr}^{0.5}$. Uncertainty margin $+15\%$ matches $+1\sigma$ on $\ln k$. \\
\addlinespace[1pt]
\textbf{Strength Model ($f_c^{\mathrm{model}}$)} & \textbf{Tier 2} & Multivariable OLS on $B/W$, $\mathrm{Air}\%$, and SCM fractions & Bolomey solid-to-void ratio with air void volume correction & Sample RMSE $3.78\ \mathrm{MPa}$, $s_e = 4.11\ \mathrm{MPa}$, LOO-CV $4.38\ \mathrm{MPa}$. Buffer $+8.25\ \mathrm{MPa}$ is IS~10262 codified margin ($+4.4\ \mathrm{MPa}$ sensitivity). \\
\addlinespace[1pt]
\textbf{Binder Domain ($B$)} & \textbf{Tier 1/3} & Batch records: SCM mixes exist only for $B \in [250, 400]\ \mathrm{kg/m^3}$ & Structural minimum binder rules & Mix C1 is sole mix at $B = 220\ \mathrm{kg/m^3}$ ($0\%$ SCM). Structural baselines enforce $B \ge 280$ (XC3) and $B \ge 300$ (XC4). \\
\addlinespace[1pt]
\textbf{50-Year Depth ($d_{50}$)} & \textbf{Tier 3} & Field monitoring terminates at $10.1\text{--}11.5\ \mathrm{yr}$ (mean 10.75\,yr) & Fickian diffusion projection ($t = 50\ \mathrm{years}$) & Extrapolates $\approx 4.3\text{--}5.0\times$ beyond field observations; subject to $+15\%$ kinetic uncertainty. \\
\addlinespace[1pt]
\textbf{$\mathrm{CO}_2$ Mass Uptake} & \textbf{Tier 3} & Measured depth $x_c$ and binder content $B$ & Stoichiometric mineralization chemistry & Cohort uptake offsets $3.8\%\text{--}5.7\%$ (baseline $4.7\%$, optimal $4.7\%\text{--}7.7\%$, benchmark $6.4\%\text{--}6.8\%$, $9.1\text{--}9.8\ \mathrm{kg/m^3}$). \\
\addlinespace[1pt]
\textbf{Embodied Carbon ($GWP$)} & \textbf{Tier 3} & Mix constituent masses in $\mathrm{kg/m^3}$ & Standard LCI emission factors (ICE v3.0 \cite{ice3}, Scrivener et al. \cite{scrivener2018}) & Cradle-to-gate A1--A3 minus Module B1 uptake. Baseline $0.840\ \mathrm{kg/kg}$ OPC; Indian plants ($0.893\text{--}0.940$ \cite{karadumpa2024}) scale savings. \\
\bottomrule
\end{tabularx}
\end{table*}

\begin{figure*}[!t]
\centering
\safefig[width=0.80\textwidth]{fig4_strength_verification_abrams_air.png}
\caption{Compressive strength modeling, Bolomey relation verification, and air void normalization across all $N = 45$ concrete mixtures: (a) Apparent breakdown of classical strength vs. $B/W$ relationship prior to air void correction ($R^2 = 0.071$ against $B/W$, or $R^2 = 0.095$ against $w/b$); (b) Air-normalized compressive strength ($f_{c,\mathrm{norm}}$ normalized to $2.0\%$ baseline air content) vs. $B/W$, restoring the fundamental linear Bolomey trend across supplementary cementitious material replacement levels ($0\%\text{--}50\%$).}
\label{fig:strength_verification}
\end{figure*}

\section{Predictive Modeling: Strength, Kinetics, and Stoichiometry}

\subsection{Gel-Space Void Ratio and Compressive Strength Formulation}
Compressive strength in hardened cementitious systems is dictated by paste porosity and solid-to-void ratio. Following Powers' gel-space ratio $\mathcal{X} = V_{\mathrm{gel}} / (V_{\mathrm{gel}} + V_{\mathrm{cap}} + V_{\mathrm{air}})$ and Feret's law:
\begin{align}
f_c &= K_F \left( \frac{V_B}{V_B + V_W + V_{\mathrm{air}}} \right)^2 \nonumber \\
&\approx K_B \left( \frac{B/\rho_B}{B/\rho_B + W/\rho_W + V_{\mathrm{air}}} - c_0 \right)
\label{eq:feret_bolomey}
\end{align}
Linearizing Eq.~\ref{eq:feret_bolomey} with respect to $B/W = 1/(w/b)$ and entrapped air volume $\mathrm{Air}\%$ yields the multivariable Bolomey formulation:
\begin{align}
f_{c,28}(\mathbf{x}) ={}& \beta_0 + \beta_{B/W}\left(\frac{B}{W}\right) + \beta_{\mathrm{air}}(\mathrm{Air}\%) \nonumber \\
&{}+ \beta_{\mathrm{FA}} p_{\mathrm{FA}} + \beta_{\mathrm{SG}} p_{\mathrm{SG}} + \beta_{\mathrm{LS}} p_{\mathrm{LS}} + \beta_{\mathrm{PO}} p_{\mathrm{PO}}
\label{eq:strength_model}
\end{align}
where $p_j$ are SCM mass percentages ($0\%\text{--}50\%$). Uncorrected regression of $f_c$ against $B/W$ shows breakdown ($R^2 = 0.071$, Fig.~\ref{fig:strength_verification}a), as low-$w/b$ batches had high air ($3.9\%\text{--}6.6\%$, mean $5.1\%$; $w/b$--air $r = -0.585$, $B/W$--air $r = 0.642$; VIFs $\le 2.09$, ruling out collinearity). Normalizing for batch air volume restores the linear Bolomey trend ($R^2 = 0.734$, adjusted $R^2 = 0.692$, Fig.~\ref{fig:strength_verification}b). The empirical air penalty ($\beta_{\mathrm{air}} = -4.354\ \mathrm{MPa/\%}$, $p < 10^{-6}$) reflects compaction void spaces. In practice, achieving $\le 2.0\%$ air at $w/b = 0.45$ ($126\ \mathrm{L/m^3}$ water) relies on polycarboxylate ether superplasticizers with defoamers (Appendix~\ref{sec:appendix_standards}).

\subsection{\texorpdfstring{Model Calibration, Leverage, and $+8.25\ \mathrm{MPa}$ Design Buffer}{Model Calibration, Leverage, and +8.25 MPa Design Buffer}}
Calibration of Eq.~\ref{eq:strength_model} across all $N = 45$ concrete mixtures gives the parameter estimates in Table~\ref{tab:strength_models}. Slag exhibits latent hydraulic reactivity with a modest slope ($\beta_{\mathrm{SG}} = -0.098\ \mathrm{MPa/\%}$, $p = 0.107$, 95\% CI: $[-0.218, +0.022]$), whereas fly ash incurs a severe early penalty ($\beta_{\mathrm{FA}} = -0.421\ \mathrm{MPa/\%}$, $p < 10^{-5}$).

\begin{table}[!t]
\caption{Air-Normalized Bolomey Compressive Strength Model Parameters and Error Metrics ($N = 45$)}
\label{tab:strength_models}
\centering
\scriptsize
\renewcommand{\arraystretch}{0.70}
\begin{tabularx}{\columnwidth}{>{\raggedright\arraybackslash}X C{2.6cm}}
\toprule
\textbf{Model Parameter / Metric} & \textbf{Estimate / Value} \\
\midrule
Intercept ($\beta_0$) & $-3.698\ \mathrm{MPa}$ \\
Bolomey Slope ($\beta_{B/W}$) & $+30.959\ \mathrm{MPa}$ \\
Air Void Penalty ($\beta_{\mathrm{air}}$) & $-4.354\ \mathrm{MPa/\%}$ \\
Fly Ash Slope ($\beta_{\mathrm{FA}}$) & $-0.421\ \mathrm{MPa/\%}$ \\
Slag Slope ($\beta_{\mathrm{SG}}$) & $-0.098\ \mathrm{MPa/\%}$ \\
Limestone Slope ($\beta_{\mathrm{LS}}$) & $-0.467\ \mathrm{MPa/\%}$ \\
Pozzolan Slope ($\beta_{\mathrm{PO}}$) & $-0.589\ \mathrm{MPa/\%}$ \\
\midrule
Coefficient of Determination ($R^2$ / $\text{Adj } R^2$) & $0.734$ / $0.692$ \\
Sample RMSE ($N = 45$) & $\mathbf{3.78\ \mathrm{MPa}}$ \\
Residual Standard Error ($s_e$, $\text{df} = 38$) & $\mathbf{4.11\ \mathrm{MPa}}$ \\
Leave-One-Out CV Error ($\mathrm{RMSE}_{\mathrm{CV}}$) & $\mathbf{4.38\ \mathrm{MPa}}$ ($R^2_{\mathrm{CV}} = 0.644$) \\
Codified Screening Buffer ($\Delta f_{\mathrm{margin}}$, IS~10262) & $\mathbf{+8.25\ \mathrm{MPa}}$ \\
\bottomrule
\end{tabularx}
\end{table}

Table~\ref{tab:strength_models} distinguishes sample RMSE ($\text{df}=45$), residual error ($\text{df}=38$), and LOO-CV error computed via diagonal hat-matrix leverage $h_{ii}$. For any out-of-sample design mixture $\mathbf{x}_0$, prediction variance accounts for estimation uncertainty through leverage $h_{00} = \mathbf{x}_0^T (\mathbf{X}^T \mathbf{X})^{-1} \mathbf{x}_0$:
\begin{equation}
\begin{aligned}
s_{\mathrm{pred}}^2(\mathbf{x}_0)
&= s_e^2 \left( 1 + \mathbf{x}_0^T (\mathbf{X}^T \mathbf{X})^{-1} \mathbf{x}_0 \right) \\
&= s_e^2 (1 + h_{00})
\end{aligned}
\label{eq:pred_variance}
\end{equation}
For boundary corners ($h_{00} \approx 0.250 > \bar{h} = 0.156$), $s_{\mathrm{pred}} = s_e\sqrt{1 + h_{00}} = 4.60\ \mathrm{MPa}$ (exact $4.6007\ \mathrm{MPa}$). We enforce a \textbf{codified design screening buffer}:
\begin{equation}
\begin{aligned}
f_{c,28}^{\mathrm{model}}(\mathbf{x})
&\ge f_{\mathrm{target}} + \Delta f_{\mathrm{margin}} \\
&= f_{\mathrm{target}} + 8.25\ \mathrm{MPa}
\end{aligned}
\label{eq:strength_buffer}
\end{equation}
from IS~10262:2019 Table 2 ($s=5.0\ \mathrm{MPa}$ for M30--M60; $+1.65 s = +8.25\ \mathrm{MPa}$ exceeds Table 1's $X = 6.5\ \mathrm{MPa}$). Sensitivities retain $s=4.0\ \mathrm{MPa}$ ($+6.6\ \mathrm{MPa}$) and LOO-CV $+4.4\ \mathrm{MPa}$. Target demands $f_{\mathrm{target}} = 30, 35, 40\ \mathrm{MPa}$ represent cylinder characteristic strengths $f_{ck,\mathrm{cyl}}$ (tested per EN~12390-3), mapping conservatively to cube grades M40, M45, and M50 per EN~206 and IS~456 ($f_{ck,\mathrm{cube}} \approx 1.25 f_{ck,\mathrm{cyl}}$). This yields continuous design thresholds of \textbf{$38.25\ \mathrm{MPa}$} for C30/37, \textbf{$43.25\ \mathrm{MPa}$} for C35/45, and \textbf{$48.25\ \mathrm{MPa}$} for C40/50. For discrete mixtures, nominal screening uses measured $f_{c,\mathrm{meas}} \ge f_{\mathrm{target}}$; buffered screening enforces $f_{c,\mathrm{meas}} \ge f_{\mathrm{target}} + 8.25\ \mathrm{MPa}$.

\subsection{Carbonation Kinetic Modeling and SCM Contrast Ratio}
Natural carbonation follows Fickian diffusion: $x_c(t) = k\sqrt{t}$, where $k$ is regressed through the origin across the six epochs: $k = \sum (x_{c,i}\sqrt{t_i}) / \sum t_i$ (Fig.~\ref{fig:kinetics_condition}). Multivariable log-linear models capture binder chemistry, $w/b$, and microclimate ($q_j = \mathrm{SCM}_j\% / 10.0$):
\begin{align}
\ln k_{\mathrm{sh}} ={}& -1.2939 + 4.0681(w/b) + 0.2525 q_{\mathrm{FA}} \nonumber \\
&{}+ 0.1161 q_{\mathrm{SG}} + 0.3089 q_{\mathrm{LS}} + 0.2555 q_{\mathrm{PO}}
\label{eq:kin_sh}
\end{align}
($R^2 = 0.876,\; F = 54.92,\; p < 10^{-15}$ for Sheltered XC3), and:
\begin{align}
\ln k_{\mathrm{un}} ={}& -1.8130 + 4.4797(w/b) + 0.2482 q_{\mathrm{FA}} \nonumber \\
&{}+ 0.0986 q_{\mathrm{SG}} + 0.3135 q_{\mathrm{LS}} + 0.2543 q_{\mathrm{PO}}
\label{eq:kin_un}
\end{align}
($R^2 = 0.863,\; F = 49.15,\; p < 10^{-15}$ for Unsheltered XC4). In log space, fly ash exhibits a marginal kinetic slope $\gamma_{\mathrm{FA}} / \gamma_{\mathrm{SG}} = 0.2525 / 0.1161 = \mathbf{2.17\times}$ steeper than slag: $+10\%$ fly ash accelerates carbonation by $+28.7\%$ ($\exp(0.2525) = 1.287$), whereas $+10\%$ slag increases $k$ by only $+12.3\%$ ($\exp(0.1161) = 1.123$).

For sheltered XC3, in-sample $\mathrm{RMSE} = 0.1411$, $s_e = 0.1516$ ($\text{df} = 39$), and $\mathrm{RMSE}_{\mathrm{CV}} = 0.1663$; for unsheltered XC4, $\mathrm{RMSE} = 0.1552$, $s_e = 0.1667$, and $\mathrm{RMSE}_{\mathrm{CV}} = 0.1765$. As an operational benchmark, we enforce a $+1\sigma$ margin ($\exp(0.1411) = 1.1515 \approx 1.15$):
\begin{equation}
k_{\mathrm{design}} = \exp(\ln k_{\mathrm{model}} + 0.1411) \approx 1.15\,k_{\mathrm{model}}
\label{eq:k_margin}
\end{equation}
imposing a \textbf{$+15\%$ kinetic uncertainty margin} (one-sided $+1\sigma$ band covering $\approx 84\%$ of normal cases). For 50-year depth $d_{50} = k_{\mathrm{model}}\sqrt{50}$, enforcing $d_{50,\mathrm{design}} \le d_{\mathrm{cover}}$ requires:
\begin{align}
d_{50,\mathrm{design}} &= 1.15 \, k_{\mathrm{model}}\sqrt{50} \le d_{\mathrm{cover}} \nonumber \\
\iff d_{50} &\le \frac{d_{\mathrm{cover}}}{1.15} = d_{50,\mathrm{allow}} = 0.8696 \cdot d_{\mathrm{cover}}
\label{eq:d50_allow}
\end{align}
yielding $d_{50,\mathrm{allow}} = 21.74\ \mathrm{mm}$ for $25\ \mathrm{mm}$ cover, $26.09\ \mathrm{mm}$ for $30\ \mathrm{mm}$, and $30.43\ \mathrm{mm}$ for $35\ \mathrm{mm}$ cover.

\subsection{\texorpdfstring{Stoichiometric $\mathrm{CO}_2$ Mineralization and Chebyshev's Sum Inequality}{Stoichiometric CO2 Mineralization and Chebyshev's Sum Inequality}}
Atmospheric carbonation mineralizes calcium hydroxide and C-S-H:
\begin{align}
\mathrm{Ca(OH)}_2 + \mathrm{CO}_2 &\to \mathrm{CaCO}_3 + \mathrm{H}_2\mathrm{O} \label{eq:rxn1} \\
\mathrm{C}\text{-}\mathrm{S}\text{-}\mathrm{H} + 2\,\mathrm{CO}_2 &\to \mathrm{SiO}_2 \cdot m\mathrm{H}_2\mathrm{O} + 2\,\mathrm{CaCO}_3 \label{eq:rxn2}
\end{align}
Following standardized uptake accounting (EN 16757, Module B1; Appendix~\ref{sec:appendix_standards}) and Lagerblad \cite{lagerblad2005}, volumetric use-phase $\mathrm{CO}_2$ uptake ($\mathrm{kg\ CO_2 / m^3}$) over 50 years is:
\begin{align}
\mathrm{CO}_{2,\mathrm{uptake}}^{\mathrm{B1}}(50\mathrm{y}) ={}& x_c(50\mathrm{y}) \cdot 10^{-3} \cdot B \cdot \left(\frac{S}{V}\right) \nonumber \\
&{}\cdot \Phi_{\mathrm{chem}} \left(\frac{44.01}{56.08}\right)
\label{eq:uptake}
\end{align}
where $S/V = 5.0\ \mathrm{m^2/m^3}$ (standard 200\,mm structural slab exposed on one face) and $\Phi_{\mathrm{chem}} = \sum_i f_i \cdot \mathrm{CaO}_i \cdot \alpha_{c,i}$ (Table~\ref{tab:uptake_scenarios}). For the C35/45 benchmark defined in Section V-D, evaluating Eq.~\ref{eq:uptake} using constituent products ($\Phi_{\mathrm{chem}} = 0.296$) gives $9.75\ \mathrm{kg\ CO_2/m^3}$ ($6.8\%$ offset of $142.5\ \mathrm{kg/m^3}$ embodied GWP, Net GWP $= 132.8\ \mathrm{kg/m^3}$). Decoupled evaluation ($\overline{\mathrm{CaO}} \cdot \bar{\alpha}_c = 0.441 \times 0.626 = 0.276$) gives $9.10\ \mathrm{kg\ CO_2/m^3}$ ($6.4\%$ offset, Net GWP $= 133.4\ \mathrm{kg/m^3}$).

\begin{table}[!t]
\caption{Stoichiometric Parameters for Bounding $\mathrm{CO}_2$ Uptake Scenarios}
\label{tab:uptake_scenarios}
\centering
\scriptsize
\renewcommand{\arraystretch}{0.70}
\begin{tabularx}{\columnwidth}{L{2.4cm} c Y Y Y}
\toprule
\textbf{Binder Constituent} & \textbf{Param.} & \textbf{Low} & \textbf{Mid (Base)} & \textbf{High} \\
\midrule
\multirow{2}{*}{\textbf{OPC Clinker}} & $\mathrm{CaO}_{\mathrm{OPC}}$ & $0.62$ & $\mathbf{0.65}$ & $0.67$ \\
 & $\alpha_{c,\mathrm{OPC}}$ & $0.60$ & $\mathbf{0.70}$ & $0.80$ \\
\addlinespace[1pt]
\multirow{2}{*}{\textbf{Slag (GGBS)}} & $\mathrm{CaO}_{\mathrm{SG}}$ & $0.36$ & $\mathbf{0.40}$ & $0.44$ \\
 & $\alpha_{c,\mathrm{SG}}$ & $0.45$ & $\mathbf{0.60}$ & $0.70$ \\
\addlinespace[1pt]
\multirow{2}{*}{\textbf{Fly Ash (FA)}} & $\mathrm{CaO}_{\mathrm{FA}}$ & $0.02$ & $\mathbf{0.05}$ & $0.08$ \\
 & $\alpha_{c,\mathrm{FA}}$ & $0.35$ & $\mathbf{0.50}$ & $0.55$ \\
\addlinespace[1pt]
\multirow{2}{*}{\textbf{Limestone (LS)}} & $\mathrm{CaO}_{\mathrm{LS}}$ & $0.00$ & $\mathbf{0.00}$ & $0.00$ \\
 & $\alpha_{c,\mathrm{LS}}$ & $0.00$ & $\mathbf{0.00}$ & $0.00$ \\
\addlinespace[1pt]
\multirow{2}{*}{\textbf{Pozzolan (PO)}} & $\mathrm{CaO}_{\mathrm{PO}}$ & $0.03$ & $\mathbf{0.05}$ & $0.08$ \\
 & $\alpha_{c,\mathrm{PO}}$ & $0.35$ & $\mathbf{0.50}$ & $0.55$ \\
\bottomrule
\end{tabularx}
\end{table}

Constituent calcium contents and carbonation reactivities have the same descending order across the primary constituents:
\begin{align}
\mathrm{CaO}_{\mathrm{OPC}} &> \mathrm{CaO}_{\mathrm{SG}} > \mathrm{CaO}_{\mathrm{FA}} \ge \mathrm{CaO}_{\mathrm{LS}} \label{eq:cheb_order1} \\
\alpha_{c,\mathrm{OPC}} &> \alpha_{c,\mathrm{SG}} > \alpha_{c,\mathrm{FA}} \ge \alpha_{c,\mathrm{LS}} \label{eq:cheb_order2}
\end{align}
Chebyshev's algebraic sum inequality guarantees that for any convex blending fractions ($f_i \ge 0, \sum f_i = 1$):
\begin{equation}
\begin{aligned}
\sum_{i=1}^m f_i \mathrm{CaO}_i \alpha_{c,i}
&\ge \left( \sum_{i=1}^m f_i \mathrm{CaO}_i \right) \left( \sum_{i=1}^m f_i \alpha_{c,i} \right) \\
&= \overline{\mathrm{CaO}} \cdot \bar{\alpha}_c
\end{aligned}
\label{eq:chebyshev}
\end{equation}
Thus, decoupled evaluation is a guaranteed lower bound within the stoichiometric model, underestimating uptake by $\approx 0.65\ \mathrm{kg/m^3}$ and yielding conservative Net GWP. Cohort uptake offsets $3.8\%\text{--}5.7\%$ of initial GWP (mean $4.7\%$, $11.4\ \mathrm{kg/m^3}$); for optimal mixes, uptake offsets $4.7\%\text{--}7.7\%$ ($6.4\%\text{--}6.8\%$ benchmark).

\begin{table*}[!t]
\caption{Optimal Mix and Net Carbon Comparison Across 18 Design Scenarios}
\label{tab:candidates_matrix}
\centering
\footnotesize
\renewcommand{\arraystretch}{1.0}
\setlength{\tabcolsep}{2.0pt}
\begin{minipage}[t]{0.49\textwidth}
\centering
\textbf{A. Sheltered (XC3)}\par\vspace{3pt}
\begin{tabular*}{\linewidth}{@{\extracolsep{\fill}}lc cc cc cc@{}}
\toprule
& & \multicolumn{2}{c}{\textbf{Nominal}} & \multicolumn{2}{c}{\textbf{Buffered}} & \multicolumn{2}{c}{\textbf{Continuous}} \\
\cmidrule(lr){3-4}\cmidrule(lr){5-6}\cmidrule(lr){7-8}
\textbf{Grade} & \textbf{Cover} & \textbf{Mix} & \textbf{GWP} & \textbf{Mix} & \textbf{GWP} & \textbf{Base} & \textbf{+15\%} \\
\midrule
C30/37 & 25 & \textbf{C39} & 169.2 & \textbf{C23} & 197.1 & \textbf{135.7} & \textbf{137.5} \\
C30/37 & 30 & \textbf{C24} & 147.6 & \textbf{C24} & 147.6 & \textbf{133.4} & \textbf{135.2} \\
C30/37 & 35 & \textbf{C24} & 147.6 & \textbf{C24} & 147.6 & \textbf{131.6} & \textbf{133.3} \\
\addlinespace[2pt]
C35/45 & 25 & \textbf{C39} & 169.2 & \textbf{C23} & 197.1 & \textbf{135.7} & \textbf{137.5} \\
C35/45 & 30 & \textbf{C24} & 147.6 & \textbf{C24} & 147.6 & \textbf{133.4} & \textbf{135.2} \\
C35/45 & 35 & \textbf{C24} & 147.6 & \textbf{C24} & 147.6 & \textbf{133.2} & \textbf{133.3} \\
\addlinespace[2pt]
C40/50 & 25 & \textbf{C39} & 169.2 & \textbf{C23} & 197.1 & \textbf{135.8} & \textbf{137.5} \\
C40/50 & 30 & \textbf{C24} & 147.6 & \textbf{C23} & 197.1 & \textbf{135.8} & \textbf{135.8} \\
C40/50 & 35 & \textbf{C24} & 147.6 & \textbf{C23} & 197.1 & \textbf{135.8} & \textbf{135.8} \\
\bottomrule
\end{tabular*}
\end{minipage}
\hfill
\begin{minipage}[t]{0.49\textwidth}
\centering
\textbf{B. Unsheltered (XC4)}\par\vspace{3pt}
\begin{tabular*}{\linewidth}{@{\extracolsep{\fill}}lc cc cc cc@{}}
\toprule
& & \multicolumn{2}{c}{\textbf{Nominal}} & \multicolumn{2}{c}{\textbf{Buffered}} & \multicolumn{2}{c}{\textbf{Continuous}} \\
\cmidrule(lr){3-4}\cmidrule(lr){5-6}\cmidrule(lr){7-8}
\textbf{Grade} & \textbf{Cover} & \textbf{Mix} & \textbf{GWP} & \textbf{Mix} & \textbf{GWP} & \textbf{Base} & \textbf{+15\%} \\
\midrule
C30/37 & 25 & \textbf{C24} & 151.2 & \textbf{C24} & 151.2 & \textbf{142.2} & \textbf{143.7} \\
C30/37 & 30 & \textbf{C24} & 151.2 & \textbf{C24} & 151.2 & \textbf{142.0} & \textbf{142.0} \\
C30/37 & 35 & \textbf{C24} & 151.2 & \textbf{C24} & 151.2 & \textbf{142.0} & \textbf{142.0} \\
\addlinespace[2pt]
C35/45 & 25 & \textbf{C24} & 151.2 & \textbf{C24} & 151.2 & \textbf{144.2} & \textbf{144.2} \\
C35/45 & 30 & \textbf{C24} & 151.2 & \textbf{C24} & 151.2 & \textbf{144.2} & \textbf{144.2} \\
C35/45 & 35 & \textbf{C24} & 151.2 & \textbf{C24} & 151.2 & \textbf{144.2} & \textbf{144.2} \\
\addlinespace[2pt]
C40/50 & 25 & \textbf{C24} & 151.2 & \textbf{C23} & 200.1 & \textbf{146.9} & \textbf{146.9} \\
C40/50 & 30 & \textbf{C24} & 151.2 & \textbf{C23} & 200.1 & \textbf{146.9} & \textbf{146.9} \\
C40/50 & 35 & \textbf{C24} & 151.2 & \textbf{C23} & 200.1 & \textbf{146.9} & \textbf{146.9} \\
\bottomrule
\end{tabular*}
\end{minipage}
\vspace{4pt}
\begin{minipage}{\textwidth}
\footnotesize
Cover: mm; GWP: kg CO$_2$e/m$^3$. Discrete nominal: measured strength, no margin; buffered: +8.25 MPa and +15\% kinetics. Continuous: +8.25 MPa strength with baseline or +15\% kinetics. Compositions, screening counts, conditions, and qualifications: Appendix~\ref{sec:appendix_candidates}.
\end{minipage}
\end{table*}

\section{Continuous Lifecycle Optimization Formulation}

\subsection{Modeling Assumptions}
\noindent\begin{minipage}{\columnwidth}
The optimization model uses six explicit assumptions:
\begin{enumerate}\setlength{\itemsep}{0.5pt}\setlength{\parskip}{0pt}
    \item[\textbf{A1}] \textbf{1D Fickian Diffusion}: Sharp mineralization front $x_c(t) = k\sqrt{t}$ \cite{papadakis1991}.
    \item[\textbf{A2}] \textbf{Decoupled Kinetics}: Rate $\ln k$ is log-linear in $w/b$ and SCM fractions ($\alpha = 0$).
    \item[\textbf{A3}] \textbf{Bolomey Strength}: 28-d strength is linear in $B/W$, $\mathrm{Air}\%$, and SCM substitution.
    \item[\textbf{A4}] \textbf{Linear LCI}: Cradle-to-gate carbon (A1--A3) is strictly linear and additive \cite{ice3}.
    \item[\textbf{A5}] \textbf{Stoichiometric Uptake}: B1 uptake is lower-bounded via Chebyshev's inequality.
    \item[\textbf{A6}] \textbf{Compaction}: Nominal air $V_{\mathrm{air}} = 2.0\%$ via superplasticizer (Appendix~\ref{sec:appendix_standards}).
\end{enumerate}
\end{minipage}

\subsection{Decision Vector, Objective, and Numbered Constraints}
The continuous mix optimization problem is defined by the decision vector 
\[
\mathbf{x} = \left[ w/b, \; B, \; x_{\mathrm{FA}}, \; x_{\mathrm{SG}}, \; x_{\mathrm{LS}}, \; x_{\mathrm{PO}} \right]^T \in \mathbb{R}^6
\]
where $x_j \in [0, 1]$ represents the mass fraction of cementitious constituent $j$. Clinker mass fraction is determined by mass conservation $x_{\mathrm{OPC}} = 1.0 - \sum_{j \in \mathrm{SCM}} x_j$. Constituent masses per cubic meter ($m_i = B \cdot x_i$) and batch water ($m_{\mathrm{w}} = B \cdot (w/b)$) satisfy absolute volume conservation:
\begin{equation}
m_{\mathrm{agg}} = \left( 1.0 - \sum_i \frac{m_i}{\rho_i} - \frac{m_{\mathrm{w}}}{\rho_{\mathrm{w}}} - V_{\mathrm{air}} \right) \cdot \rho_{\mathrm{agg}}
\label{eq:agg_balance}
\end{equation}
using the following constituent material densities (all in $\mathrm{kg/m^3}$): $\rho_{\mathrm{OPC}} = 3150$, $\rho_{\mathrm{FA}} = 2300$, $\rho_{\mathrm{SG}} = 2900$, $\rho_{\mathrm{LS}} = 2700$, $\rho_{\mathrm{PO}} = 2500$, $\rho_{\mathrm{agg}} = 2650$, and $\rho_{\mathrm{w}} = 1000$.

Initial cradle-to-gate embodied carbon (Modules A1--A3) is:
\begin{align}
GWP_{\mathrm{A1-A3}}(\mathbf{x}) ={}& \sum_i m_i EF_i + m_{\mathrm{agg}}EF_{\mathrm{agg}} \nonumber \\
&{}+ m_{\mathrm{w}}EF_{\mathrm{w}} + EF_{\mathrm{batching}}
\label{eq:gwp_a1a3}
\end{align}
where standard emission factors are $EF_{\mathrm{OPC}}=0.840$, $EF_{\mathrm{FA}}=0.015$, $EF_{\mathrm{SG}}=0.060$, $EF_{\mathrm{LS}}=0.025$, $EF_{\mathrm{PO}}=0.030$, $EF_{\mathrm{agg}}=0.006$, $EF_{\mathrm{w}}=0.0005\ \mathrm{kg\ CO_2e/kg}$ \cite{ice3}, and $EF_{\mathrm{batching}} = 7.5\ \mathrm{kg\ CO_2e/m^3}$.

The 50-year use-phase stoichiometric $\mathrm{CO}_2$ uptake (Module B1) is:
\begin{align}
\mathrm{CO}_{2,\mathrm{uptake}}^{\mathrm{B1}}(\mathbf{x}, 50\mathrm{y}) ={}& x_c(50\mathrm{y}, \mathbf{x}) \cdot 10^{-3} \cdot B \cdot \left(\frac{S}{V}\right) \nonumber \\
&{}\cdot (\overline{\mathrm{CaO}} \cdot \bar{\alpha}_c) \left(\frac{44.01}{56.08}\right)
\label{eq:uptake_formal}
\end{align}
where $S/V = 5.0\ \mathrm{m^2/m^3}$. Net lifecycle GWP subtracts estimated use-phase $\mathrm{CO}_2$ uptake from production emissions:
\begin{equation}
\begin{aligned}
\min_{\mathbf{x}} \; \mathcal{J}(\mathbf{x})
&= GWP_{\mathrm{A1-A3}}(\mathbf{x}) \\
&\quad - \mathrm{CO}_{2,\mathrm{uptake}}^{\mathrm{B1}}(\mathbf{x}, 50\mathrm{y})
\end{aligned}
\label{eq:nlp_opt}
\end{equation}
subject to the following numbered constraints:
\begin{align}
g_{\mathrm{str}}(\mathbf{x}) &= (f_{\mathrm{target}} + \Delta f_{\mathrm{margin}}) - f_{c,28}^{\mathrm{model}}(\mathbf{x}) \le 0 \label{eq:cstr_str} \\
g_{\mathrm{dur}}(\mathbf{x}) &= \gamma_{\mathrm{dur}} \, x_c(50\mathrm{y}, \mathbf{x}) - d_{\mathrm{cover}} \le 0 \label{eq:cstr_dur} \\
g_{\mathrm{scm}}(\mathbf{x}) &= \sum_{j \in \mathrm{SCM}} x_j - 0.50 \le 0 \label{eq:cstr_scm} \\
g_{\mathrm{FA}}(\mathbf{x}) &= x_{\mathrm{FA}} - 0.35 \le 0 \label{eq:cstr_fa} \\
g_{\mathrm{bnd}}(\mathbf{x}) &: w/b \le w/b_{\max}, \nonumber \\
&\quad B \ge B_{\min}, \; x_j \ge 0 \label{eq:cstr_bounds}
\end{align}
where $\gamma_{\mathrm{dur}}=1.0$ (nominal) or $1.15$ (uncertain); strength margins follow Section III-B. Bound limits enforce structural baselines ($B \ge 280\ \mathrm{kg/m^3}$ for XC3; $B \ge 300\ \mathrm{kg/m^3}$ for XC4; $w/b \le 0.55$ for XC3; $w/b \le 0.50$ for XC4).

The Lagrangian associated with Eq.~\ref{eq:nlp_opt} is $\mathcal{L}(\mathbf{x}, \boldsymbol{\mu}) = \mathcal{J}(\mathbf{x}) + \sum_k \mu_k g_k(\mathbf{x})$. At local optimum $\mathbf{x}^*$, the first-order Karush-Kuhn-Tucker (KKT) conditions require:
\begin{align}
\nabla_{\mathbf{x}} \mathcal{L}(\mathbf{x}^*, \boldsymbol{\mu}^*) &= \mathbf{0} \label{eq:kkt_stationarity} \\
&\begin{aligned}
\mu_k^* g_k(\mathbf{x}^*) = 0, \quad \mu_k^* &\ge 0, \\
g_k(\mathbf{x}^*) &\le 0 \quad \forall k
\end{aligned}\label{eq:kkt_complementarity}
\end{align}
By the envelope theorem, the dual multiplier $\mu_k^* = -\partial \mathcal{J}^*/\partial b_k$ represents the marginal shadow price of constraint $k$.
In Section~\ref{sec:results}, we link $\mu_{\mathrm{dur}}^*$ directly to the carbon dividend of structural cover.

\begin{figure*}[!t]
\centering
\safefig[width=0.80\textwidth]{fig6_pareto_carbon_vs_strength.png}
\caption{Multi-objective Pareto frontiers: (a) Discrete mix landscape with buffered strength benchmarks; (b) Continuous Pareto trade-off frontiers across covers ($25, 30, 35\ \mathrm{mm}$) contrasting unconstrained envelope ($B \ge 220\ \mathrm{kg/m^3}$) against structural baseline ($B \ge 280\ \mathrm{kg/m^3}$, XC3).}
\label{fig:pareto}
\end{figure*}

\noindent\textbf{Structural Baselines vs.\ Extrapolation:} Codified binder and FA baselines enforce minimum cover ($d_{50} = d_{\mathrm{cover}}$). Section VI discusses paste starvation \cite{neville2011} and unvalidated SCM extrapolation at $B=220\ \mathrm{kg/m^3}$.

\section{Lifecycle Optimization Results and Performance Analysis}
\label{sec:results}

\subsection{Performance Across the Eighteen Design Scenarios}
Table~\ref{tab:candidates_matrix} summarizes optimal concrete mixtures across all nine structural design cases evaluated under sheltered and unsheltered microclimates (18 total scenarios). Each case combines a target strength with a cover depth, allowing the effects of strength demand, cover, and exposure to be compared within the same framework. The table reports discrete screening alongside continuous optimization, with nominal and buffered results shown separately. Discrete winners are selected from measured mixtures, whereas continuous solutions are model-predicted compositions subject to the specified strength, durability, and binder constraints. Appendix~\ref{sec:appendix_candidates} provides the corresponding compositions and screening counts. The following subsections examine the constraints that change the selected mixtures and the carbon savings associated with continuous design.

\subsection{\texorpdfstring{Discrete Screening and the 25\,mm Uncertainty Penalty}{Discrete Screening and the 25 mm Uncertainty Penalty}}
Table~\ref{tab:candidates_matrix} gives the discrete winners; their controlling boundaries are:
\begin{enumerate}\setlength{\itemsep}{0.5pt}\setlength{\parskip}{0pt}
    \item \textbf{The 25\,mm Sheltered Dilemma (Mix C39 vs. C23)}: Nominal C39 clears C40/50 by only $0.4\ \mathrm{MPa}$. Codified margins disqualify it for $\ge\mathrm{C35/45}$, while $+15\%$ kinetics ($1.15\times22.6=26.0>25\ \mathrm{mm}$) eliminates it across all grades. C23 is the unique buffered survivor, imposing a $+16.5\%$ carbon penalty.
    \item \textbf{Dominance of Mix C24 at $\ge 30\ \mathrm{mm}$ Cover}: C24 dominates C30/37 and C35/45 under nominal, model-buffered and IS~10262 $s=4$ screens ($d_{50}=25.0\ \mathrm{mm}$, $28.8\ \mathrm{mm}$ buffered), saving $25.1\%$ ($49.5\ \mathrm{kg/m^3}$) over C23. EN~206 permits it; IS~456 Severe excludes it. The main $s=5.0\ \mathrm{MPa}$ margin eliminates C24 for C40/50 and promotes C23.
    \item \textbf{Unsheltered Microclimate Performance}: Rain pore-blocking retards C24 carbonation to $d_{50}=17.52\ \mathrm{mm}$. C24 wins C30/37 and C35/45 across covers; C23 wins C40/50 under $s=5.0\ \mathrm{MPa}$. The $w/b\le0.50$ pool spans 24/26 nominally and 5/26 for C40/50 under joint buffering.
\end{enumerate}

\begin{table*}[!t]
\caption{Systematic Sensitivity and Robustness Matrix Across Eight Optimization Dimensions for Benchmark Structural Scenario (Grade C35/45, 30\,mm Cover, Sheltered Exposure)}
\label{tab:sensitivity_matrix}
\centering
\scriptsize
\renewcommand{\arraystretch}{0.48}
\setlength{\tabcolsep}{2.5pt}
\begin{tabularx}{\textwidth}{L{2.3cm} L{2.5cm} L{5.2cm} C{1.6cm} >{\raggedright\arraybackslash}X}
\toprule
\textbf{Dimension} & \textbf{Scenario} & \textbf{Optimal Mixture ($w/b, B$, SCM\%)} & \textbf{Net GWP (kg/m$^3$)} & \textbf{Active Constraint} \\
\midrule
\multirow{2}{*}{\textbf{Exposure Class*}} & Sheltered Baseline (XC3) & $0.45, 280\ \mathrm{kg}, 24.0\%\ \mathrm{FA}, 26.0\%\ \mathrm{SG}$ & 133.4 & Durability ($d_{50} = 30.0\ \mathrm{mm}$) \\
 & Unsheltered Code (XC4) & $0.45, 300\ \mathrm{kg}, 25.5\%\ \mathrm{FA}, 24.5\%\ \mathrm{SG}$ & 144.2 & Strength ($f_c = 43.3\ \mathrm{MPa}$) \\
\addlinespace[0.5pt]
\multirow{3}{*}{\textbf{SCM Allocation}} & Economic (Baseline) & $0.45, 280\ \mathrm{kg}, 24.0\%\ \mathrm{FA}, 26.0\%\ \mathrm{SG}$ & 133.4 & Durability ($d_{50} = 30.0\ \mathrm{mm}$) \\
 & Cut-Off (Waste) & $0.51, 280\ \mathrm{kg}, 0.0\%\ \mathrm{FA}, 50.0\%\ \mathrm{SG}$ & 129.4 & Strength ($f_c = 43.3\ \mathrm{MPa}$) \\
 & Mass Allocation & $0.45, 280\ \mathrm{kg}, 24.0\%\ \mathrm{FA}, 26.0\%\ \mathrm{SG}$ & 138.8 & Durability ($d_{50} = 30.0\ \mathrm{mm}$) \\
\addlinespace[0.5pt]
\multirow{3}{*}{\textbf{Kinetic Model}} & Fickian ($n = 0.50$) & $0.45, 280\ \mathrm{kg}, 24.0\%\ \mathrm{FA}, 26.0\%\ \mathrm{SG}$ & 133.4 & Durability ($d_{50} = 30.0\ \mathrm{mm}$) \\
 & Power-Law ($n = 0.509$) & $0.45, 280\ \mathrm{kg}, 23.0\%\ \mathrm{FA}, 27.0\%\ \mathrm{SG}$ & 133.5 & Durability ($d_{50} = 30.0\ \mathrm{mm}$) \\
 & Power-Law ($n = 0.452$) & $0.45, 280\ \mathrm{kg}, 25.5\%\ \mathrm{FA}, 24.5\%\ \mathrm{SG}$ & 133.9 & Strength ($f_c = 43.3\ \mathrm{MPa}$) \\
\addlinespace[0.5pt]
\multirow{3}{*}{\textbf{$\mathrm{CO}_2$ Uptake}} & Low Uptake & $0.45, 280\ \mathrm{kg}, 24.0\%\ \mathrm{FA}, 26.0\%\ \mathrm{SG}$ & 135.8 & Durability ($d_{50} = 30.0\ \mathrm{mm}$) \\
 & Midpoint Uptake & $0.45, 280\ \mathrm{kg}, 24.0\%\ \mathrm{FA}, 26.0\%\ \mathrm{SG}$ & 133.4 & Durability ($d_{50} = 30.0\ \mathrm{mm}$) \\
 & High Uptake & $0.45, 280\ \mathrm{kg}, 24.0\%\ \mathrm{FA}, 26.0\%\ \mathrm{SG}$ & 131.5 & Durability ($d_{50} = 30.0\ \mathrm{mm}$) \\
\addlinespace[0.5pt]
\multirow{3}{*}{\textbf{Compaction Air}} & Air = 1.5\% & $0.45, 280\ \mathrm{kg}, 24.0\%\ \mathrm{FA}, 26.0\%\ \mathrm{SG}$ & 133.4 & Durability ($d_{50} = 30.0\ \mathrm{mm}$) \\
 & Air = 2.0\% (Baseline) & $0.45, 280\ \mathrm{kg}, 24.0\%\ \mathrm{FA}, 26.0\%\ \mathrm{SG}$ & 133.4 & Durability ($d_{50} = 30.0\ \mathrm{mm}$) \\
 & Air = 3.0\% & $0.45, 280\ \mathrm{kg}, 12.1\%\ \mathrm{FA}, 37.9\%\ \mathrm{SG}$ & 135.5 & Strength ($f_c = 43.3\ \mathrm{MPa}$) \\
\addlinespace[0.5pt]
\multirow{2}{*}{\textbf{SCM Rules}} & Empirical (SCM $\le 50\%$) & $0.45, 280\ \mathrm{kg}, 24.0\%\ \mathrm{FA}, 26.0\%\ \mathrm{SG}$ & 133.4 & Durability ($d_{50} = 30.0\ \mathrm{mm}$) \\
 & Code (FA $\le 35\%$) & $0.45, 280\ \mathrm{kg}, 24.0\%\ \mathrm{FA}, 26.0\%\ \mathrm{SG}$ & 133.4 & Durability ($d_{50} = 30.0\ \mathrm{mm}$) \\
\addlinespace[0.5pt]
\multirow{2}{*}{\textbf{Strength Margin}} & IS~10262 Table 2 ($s=5$, $+8.25\ \mathrm{MPa}$, Main) & $0.45, 280\ \mathrm{kg}, 24.0\%\ \mathrm{FA}, 26.0\%\ \mathrm{SG}$ & 133.4 & Durability ($0.48\ \mathrm{MPa}$ slack) \\
 & Sensitivities: $s=4$ (+6.6), LOO-CV (+4.4), Nom (0) & $0.45, 280\ \mathrm{kg}, 24.0\%\ \mathrm{FA}, 26.0\%\ \mathrm{SG}$ & 133.4 & Durability (slacks $2.13, 4.33, 8.73\ \mathrm{MPa}$) \\
\addlinespace[0.5pt]
\multirow{3}{*}{\textbf{Codified Baseline}} & EN~206 Baseline ($B \ge 280\ \mathrm{kg/m^3}$) & $0.45, 280\ \mathrm{kg}, 24.0\%\ \mathrm{FA}, 26.0\%\ \mathrm{SG}$ & 133.4 & Durability ($d_{50} = 30.0\ \mathrm{mm}$) \\
 & IS~456 Mild ($B \ge 300\ \mathrm{kg/m^3}$) & $0.45, 300\ \mathrm{kg}, 24.0\%\ \mathrm{FA}, 26.0\%\ \mathrm{SG}$ & 141.3 & Durability ($d_{50} = 30.0\ \mathrm{mm}$) \\
 & IS~456 Severe ($B \ge 320\ \mathrm{kg/m^3}$) & $0.45, 320\ \mathrm{kg}, 24.0\%\ \mathrm{FA}, 26.0\%\ \mathrm{SG}$ & 149.2 & Durability ($d_{50} = 30.0\ \mathrm{mm}$) \\
\bottomrule
\end{tabularx}
\vspace{1pt}
\begin{minipage}{\textwidth}
\tiny
*Outdoor difference ($+10.8\ \mathrm{kg/m^3}$) reflects codified binder ($B=300$ vs $280\ \mathrm{kg/m^3}$) and kinetics. IS~456 Mild ($B \ge 300\ \mathrm{kg/m^3}$) yields $141.3\ \mathrm{kg/m^3}$ ($6.3\ \mathrm{kg/m^3}$ saving over C24); Severe ($B \ge 320\ \mathrm{kg/m^3}$) gives $149.2\ \mathrm{kg/m^3}$ (C24 excluded). At $1.645\sigma$ ($7.20\ \mathrm{MPa}, +26.1\%$ kinetic), C24 reaches $31.6\ \mathrm{mm} > 30\ \mathrm{mm}$ (fails); optimum ($136.3\ \mathrm{kg/m^3}$) saves $60.8\ \mathrm{kg/m^3}$ ($30.8\%$) over survivor C23 ($197.1\ \mathrm{kg/m^3}$). 1.0\% PCE gives Net GWP $137.6\text{--}138.7\ \mathrm{kg/m^3}$ ($8.8\text{--}9.9\ \mathrm{kg/m^3}$ headline saving over un-admixed C24; $14.6\ \mathrm{kg/m^3}$ hypothetical like-for-like).
\end{minipage}
\end{table*}

\subsection{Continuous Optimization: Binder Reduction vs.\ SCM Blending}
Continuous optimization reduces binder content from $B=310$ to $280\ \mathrm{kg/m^3}$ while lowering $w/b$ from $0.50$ to $0.45$ and adjusting the SCM blend. The decomposition assigns $70.8\%\text{--}87.1\%$ of the savings to binder reduction, or approximately $79\%$ overall, and $12.9\%\text{--}29.2\%$ to SCM blending. The bilinear interaction is not statistically significant ($p > 0.80$). Tables~\ref{tab:candidates_matrix} and~\ref{tab:candidates_details} give all compositions and carbon values; the additional case-specific predictions are:
\begin{enumerate}\setlength{\itemsep}{0pt}\setlength{\parskip}{0pt}
    \item \textbf{Tight Cover (25\,mm, Sheltered)}: The nominal optimum gives $f_c=48.0\ \mathrm{MPa}$ and $d_{50}=25.0\ \mathrm{mm}$, saving $19.8\%$ over nominal C39. Under $+15\%$ uncertainty ($d_{50}\le21.74\ \mathrm{mm}$), $f_c=51.4\ \mathrm{MPa}$ and $d_{50}=21.7\ \mathrm{mm}$; saving over buffered C23 is $30.2\%$ ($27.5\%\text{--}28.1\%$ with $1.0\%$ PCE).
    \item \textbf{Standard Cover (30\,mm, Sheltered)}: Nominal predictions are $f_c=43.7\ \mathrm{MPa}$ and $d_{50}=30.0\ \mathrm{mm}$; under uncertainty ($d_{50}\le26.09\ \mathrm{mm}$), they become $f_c=47.0\ \mathrm{MPa}$ and $d_{50}=26.1\ \mathrm{mm}$.
    \item \textbf{Deep Cover (35\,mm, Sheltered)}: The nominal capped mixture reaches $d_{50}=34.8\ \mathrm{mm}$; unconstrained $35.3\%$ FA + $14.7\%$ SG reaches $35.0\ \mathrm{mm}$ with identical GWP. Under uncertainty, $d_{50}=30.4\ \mathrm{mm}$.
    \item \textbf{High Strength (30\,mm, C40/50)}: The codified $48.25\ \mathrm{MPa}$ requirement selects the mixture in Table~\ref{tab:candidates_matrix} ($f_c=48.2\ \mathrm{MPa}$, $d_{50}=24.8\ \mathrm{mm}$); $+4.4\ \mathrm{MPa}$ LOO-CV instead gives $22.0\%$ FA + $28.0\%$ SG ($133.8\ \mathrm{kg/m^3}$).
\end{enumerate}

\subsection{Closed-Form Durability Limits and Kuhn-Tucker Optimality}
For sheltered C35/45 at 30\,mm cover, nominal ($35.0\ \mathrm{MPa}$), LOO-CV ($39.4\ \mathrm{MPa}$), and codified ($43.25\ \mathrm{MPa}$) thresholds select the same mixture in Table~\ref{tab:candidates_matrix}, whose predicted strength is $43.73\ \mathrm{MPa}$. The propositions explain this invariance on the fixed economic-allocation FA--SG branch; derivations appear in Appendix~\ref{sec:appendix_propositions}.

\noindent\textbf{Proposition 1 (Closed-Form Durability Limit):}
At $B=280\ \mathrm{kg/m^3}$, $w/b=0.45$, and $p_{\mathrm{SG}}=50-p_{\mathrm{FA}}$, sheltered kinetics reduce to $\ln k_{\mathrm{sh}}=1.1172+0.01364p_{\mathrm{FA}}$. Thus, $d_{50}\le d_{\mathrm{allow}}$ imposes
\begin{equation}
p_{\mathrm{FA}}\le p_{\mathrm{dur}}
=\frac{\ln(d_{\mathrm{allow}}/\sqrt{50})-1.1172}{0.01364}.
\label{eq:prop_durability_bound}
\end{equation}
Strength and the PPC fly ash cap may impose tighter upper bounds; Appendix~\ref{sec:appendix_propositions} evaluates the benchmark.

\noindent\textbf{Proposition 2 (Conditional Boundary Optimality):}
If the full lifecycle objective decreases throughout the feasible interval on this branch, its minimum occurs at the largest admissible FA content:
\begin{equation}
\begin{gathered}
\frac{\mathrm d\mathcal J}{\mathrm dp_{\mathrm{FA}}}<0
\ \Longrightarrow\ \\
p_{\mathrm{FA}}^*=\min\{50,p_{\mathrm{dur}},p_{\mathrm{cap}},p_{\mathrm{str}}\},
\end{gathered}
\label{eq:prop_boundary_optimum}
\end{equation}
provided the interval is nonempty; $p_{\mathrm{cap}}=35\%$ and $p_{\mathrm{str}}$ are the PPC and strength bounds. Numerical derivatives support decreasing carbon on this economic-allocation branch, where benchmark durability binds and strength and FA caps remain slack. This conditional, fixed-branch result is not a global optimality proof. Cut-Off allocation instead favors $w/b=0.53$ and $0\%$ FA because low slag emissions and greater uptake alter the trade-off.

\noindent\textbf{Local Cover Dividend:}
With durability in millimetres, the local 30\,mm shadow price is $\mu_{\mathrm{dur}}^*=0.4087\approx0.41\ \mathrm{kg\ CO_2e/(m^3\cdot mm)}$, while $\mu_{\mathrm{str}}^*=0$. For C35/45 under +8.25\,MPa, strength binds before 35\,mm ($25.5\%$ FA, $133.2\ \mathrm{kg/m^3}$). The durability-only 25--35\,mm dividend therefore does not describe that full path (Table~\ref{tab:candidates_matrix}).

\noindent\textbf{Strength Slack and Compaction Limit:}
Codified strength slack is $0.48\ \mathrm{MPa}$, exhausted at $2.11\%$ air (LOO-CV: $2.99\%$). Batch air ($5.1\%$) incurs a $13.5\ \mathrm{MPa}$ penalty and fails all three thresholds. Appendix~\ref{sec:appendix_propositions} retains the baseline-air calculation.

\begin{figure*}[!t]
\centering
\safefig[width=0.80\textwidth]{fig7_sensitivity_analysis.png}
\caption{Sensitivity sweep for benchmark C35/45 (30\,mm cover): (a) Allocation ($\rho > 0.997$); (b) Kinetics and uptake; (c) Air and strength buffer; and (d) Binder minimum and PPC cap.}
\label{fig:sensitivity}
\end{figure*}

\section{Discussion: Pareto Dynamics, Sensitivity, and Structural Caveats}

\subsection{Pareto Dynamics and Sensitivity Sweeps}
Figure~\ref{fig:pareto} synthesizes the structural Pareto landscape:
\begin{enumerate}\setlength{\itemsep}{0.5pt}\setlength{\parskip}{0pt}
    \item \textbf{Discrete Binder Clusters (Fig.~\ref{fig:pareto}a)}: Pure OPC clusters at $204\text{--}354\ \mathrm{kg\ CO_2e/m^3}$; fly ash reaches $140\text{--}180\ \mathrm{kg/m^3}$ but combines strength loss with $d_{50}=45\text{--}62\ \mathrm{mm}$. Slag mixes ($30\%\text{--}50\%$ SG) dominate through moderate strength loss and buffering (Table~\ref{tab:strength_models}).
    \item \textbf{Cover Depth Carbon Dividend (Fig.~\ref{fig:pareto}b)}: Increasing cover from 25 to 35\,mm lowers the frontier by $4.1\ \mathrm{kg\ CO_2e/m^3}$ through the FA-for-SG substitution quantified in Section V-D and Appendix~\ref{sec:appendix_propositions}.
    \item \textbf{Structural Binder Penalty}: Enforcing EN~206 XC3 binder ($B \ge 280\ \mathrm{kg/m^3}$) adds $+23.3\text{--}24.1\ \mathrm{kg\ CO_2e/m^3}$ ($+21.5\%$; benchmark $+23.6\ \mathrm{kg/m^3}$) over $B = 220\ \mathrm{kg/m^3}$ to ensure adequate paste lubrication.
    \item \textbf{Optimization Gap Decomposition}: Continuous frontiers lie $14\text{--}60\ \mathrm{kg/m^3}$ below discrete batches. At sheltered C35/45, 30\,mm cover, C24 ($147.57\ \mathrm{kg/m^3}$) versus the continuous design ($133.45\ \mathrm{kg/m^3}$) gives model-predicted $\Delta GWP_{\mathrm{total}} = \mathbf{14.1\ \mathrm{kg\ CO_2e/m^3}}$, pending batch trials. Decomposition gives: Path~1 (Binder$\to$SCM) $10.0 + 4.1 = \mathbf{14.1\ \mathrm{kg/m^3}}$; Path~2 (SCM$\to$Binder) $1.8 + 12.3 = \mathbf{14.1\ \mathrm{kg/m^3}}$. Binder reduction contributes $\approx 79\%$ and SCM blending $\approx 21\%$. At 25\,mm cover, predicted savings over buffered C23 are $61.4\ \mathrm{kg/m^3}$ ($31.2\%$). The PCE-qualified comparison is retained in Table~\ref{tab:sensitivity_matrix}.
    \item \textbf{Corner and Boundary Extrapolation}: At the benchmark bounds described in Section V-D, nearest mixes in $(w/b, \mathrm{air}, \mathrm{SCM})$ space are: C28 ($d=1.09$: $w/b=0.50$, air $1.4\%$, $20\%$ FA + $20\%$ SG; $f_c$ $41.1$ vs.\ $41.7\ \mathrm{MPa}$; $k$ $4.374$ vs.\ $4.381$; $d_{50}$ $30.9$ vs.\ $31.0\ \mathrm{mm}$), C23 ($d=2.31$: $w/b=0.50$, air $1.4\%$, $30\%$ SG; $f_c$ $55.7$ vs.\ $49.2\ \mathrm{MPa}$; $k$ $2.534$ vs.\ $2.970$; $d_{50}$ $17.9$ vs.\ $21.0\ \mathrm{mm}$), and C39 ($d=2.33$: $w/b=0.50$, air $1.6\%$, $20\%$ SG + $20\%$ PO; $f_c$ $40.4$ vs.\ $37.5\ \mathrm{MPa}$; $k$ $3.193$ vs.\ $4.408$; $d_{50}$ $22.6$ vs.\ $31.2\ \mathrm{mm}$). Corner leverage $h_{00} = 0.250$ gives $s_{\mathrm{pred}} = s_e\sqrt{1 + 0.250} = 4.60\ \mathrm{MPa}$.
\end{enumerate}



Figure~\ref{fig:sensitivity} details the sensitivity sweep for benchmark C35/45 ($d_{\mathrm{cover}} = 30\ \mathrm{mm}$):
\begin{enumerate}\setlength{\itemsep}{0pt}\setlength{\parskip}{0pt}
    \item \textbf{Allocation (Fig.~\ref{fig:sensitivity}a)}: Mix rankings are invariant ($\rho > 0.997, p < 10^{-40}$) across regimes; clinker ($0.840\ \mathrm{kg/kg}$) dominates SCMs ($14\text{--}56\times$ economic: $\mathrm{EF}_{\mathrm{FA}} = 0.015, \mathrm{EF}_{\mathrm{SG}} = 0.060\ \mathrm{kg/kg}$; $7.0\times$ mass: $0.030, 0.120$; cut-off: $0.005, 0.020$).
    \item \textbf{Kinetics (Fig.~\ref{fig:sensitivity}b)}: The sheltered power-law ($n = 0.509$) gives a 50-year depth $28.91\ \mathrm{mm}$ vs.\ $26.40\ \mathrm{mm}$ Fickian ($+9.5\%$, mean $\bar{n} = 0.522$). The 10-year anchor shifts depth $+1.5\%$ ($30.44$ vs.\ $30.00\ \mathrm{mm}$), within $+15\%$ margin ($d_{50,\mathrm{allow}} = 26.1\ \mathrm{mm}$). Outdoor $n = 0.452$ gives $d_{50} = 30.0\ \mathrm{mm}$ ($29.7\%$ FA + $20.3\%$ SG).
    \item \textbf{$\mathrm{CO}_2$ Uptake (Fig.~\ref{fig:sensitivity}b)}: The uptake sweep shifts benchmark Net GWP by $+2.4 / -1.9\ \mathrm{kg/m^3}$; offset ranges are given in Section III-D and Table~\ref{tab:evidence_hierarchy}. High fly ash adds $5.8\ \mathrm{kg/m^3}$ uptake, but carbonation reaches $55\ \mathrm{mm}$ (Mix C35), depassivating $30\ \mathrm{mm}$ cover in $14.8\ \mathrm{yr}$. Clinker reduction dominates uptake.
    \item \textbf{Compaction (Fig.~\ref{fig:sensitivity}c)}: An air increase of $+1.0\%$ ($3.0\%$) incurs a $-4.35\ \mathrm{MPa}$ penalty, binding strength in Table~\ref{tab:sensitivity_matrix}. Section V-D gives the benchmark air-content failure thresholds.
    \item \textbf{Fly Ash Cap (Fig.~\ref{fig:sensitivity}d)}: The PPC cap ($\mathrm{FA} \le 35\%$, IS~1489 \cite{is1489}) does not bind indoors ($24.0\%$ FA); at 35\,mm, unconstrained FA reaches $35.3\%$ (Table~\ref{tab:candidates_matrix}), binding weakly ($35.0\%$ FA, $15.0\%$ SG, $131.6\ \mathrm{kg/m^3}$). In XC4, the cap binds: C35/45 at 30\,mm yields $142.3\ \mathrm{kg/m^3}$ ($35.0\%$ FA, $12.9\%$ SG, $2.1\%$ LS); at 25\,mm, carbon is $142.2 \to 142.3\ \mathrm{kg/m^3}$; and C30/37 at 30\,mm shifts $140.0 \to 141.0\ \mathrm{kg/m^3}$.
\end{enumerate}

\subsection{Structural Baselines and Practical Implementation}
Optima at $B = 220\ \mathrm{kg/m^3}$ present two challenges:
(1)~\textbf{Paste Starvation}: neat paste volume is $176.6\ \mathrm{L/m^3}$ ($196.6\ \mathrm{L/m^3}$ with 2\% air), below the $220\text{--}250\ \mathrm{L/m^3}$ limit \cite{neville2011}; codified baselines ($B \ge 280\ \mathrm{kg/m^3}$ sheltered, $B \ge 300\ \mathrm{kg/m^3}$ outdoor) restore paste to $224.8\ \mathrm{L/m^3}$.
(2)~\textbf{Clinker Depletion}: at $50\%$ SCM, clinker falls to $110\ \mathrm{kg/m^3}$ ($\approx 15\text{--}20\ \mathrm{kg/m^3}$ portlandite reserve \cite{taylor1997, papadakis1992}). Outdoor retardation ($26\%\text{--}30\%$) reflects pore-water blocking ($D_{\mathrm{aq}} \approx 10^{-4} D_{\mathrm{gas}}$). Storage effects were not statistically significant ($p > 0.50$). Slag mixes require 14--28\,d moist curing. Regional kiln factors ($0.893\text{--}0.940\ \mathrm{kg/kg}$ \cite{karadumpa2024}) scale savings proportionally.

\subsection{Limitations}
Six physical and jurisdictional limitations qualify these findings:
(a)~Phenolphthalein indicates pH 9--9.5, whereas steel depassivates near pH 10.5--11.5; carbonation fronts may thus precede color change.
(b)~IS~456 Table 5 rates sheltered concrete as Mild ($B \ge 300\ \mathrm{kg/m^3}, w/b \le 0.55$), where continuous optimum is $141.3\ \mathrm{kg/m^3}$ ($6.3\ \mathrm{kg/m^3}$ saving over C24; $1.8\text{--}2.9\ \mathrm{kg/m^3}$ with PCE), and coastal exposure as Severe ($B \ge 320\ \mathrm{kg/m^3}, w/b \le 0.45$), where C24 is excluded and continuous optimum is $149.2\ \mathrm{kg/m^3}$. Chloride ingress is not modelled.
(c)~Setting $d_{50} = d_{\mathrm{cover}}$ leaves zero execution tolerance ($5\text{--}10\ \mathrm{mm}$); site cover must therefore be thicker.
(d)~Curing under 7\,d increases carbonation by 25\%--50\% \cite{parrott1996}, exceeding the +15\% margin.
(e)~Strength screening uses IS~10262 Table 2 ($+8.25\ \mathrm{MPa}$ for $s=5.0\ \mathrm{MPa}$, leaving $0.48\ \mathrm{MPa}$ slack at benchmark); corner error $s_{\mathrm{pred}} = 4.60\ \mathrm{MPa}$ is within $4.13\ \mathrm{MPa}$ of nominal slack. The kinetic margin uses in-sample RMSE ($0.1411 \to +15\%$). The $1.645\sigma$ sensitivity in Table~\ref{tab:sensitivity_matrix} disqualifies C24 and compares the continuous design with feasible survivor C23.
(f)~The PCE compaction route remains uncalibrated on the physical cohort (Appendix~\ref{sec:appendix_standards}).

\noindent\textbf{Low-Binder Extrapolation Envelope:} Relaxing $B$ to $220\ \mathrm{kg/m^3}$ yields theoretical Net GWP of $108.3\text{--}111.6\ \mathrm{kg\ CO_2e/m^3}$ indoors and $108.7\text{--}112.4\ \mathrm{kg/m^3}$ outdoors ($\sim17.5\%$ reduction). The domain restriction in Section II-A and paste starvation above make these theoretical bounds rather than validated structural specifications.


\section{Conclusions and Structural Specification Guidance}
Field calibration across 11 years in Chennai supports the discrete screening and continuous design framework. Under EN~206, Mix C24 dominates sheltered cases at $\ge30\ \mathrm{mm}$ cover and outdoor C30/37--C35/45; IS~456 Severe excludes it, and the codified C40/50 strength margin selects C23. At 25\,mm sheltered cover, kinetic uncertainty replaces nominal C39 with buffered C23 and increases carbon by $16.5\%$ (Table~\ref{tab:candidates_matrix}).

Across 18 scenarios, continuous optimization predicts $4.6\%\text{--}31.2\%$ carbon savings ($1.1\%\text{--}29.0\%$ with $1.0\%$ PCE). Binder reduction contributes most of the improvement; SCM blending provides additional savings subject to the strength, durability, and composition limits. The decomposition in Section VI-A remains model-predicted and requires batch trials. Codified binder baselines restore paste lubrication, while the low-binder envelope remains unvalidated.

Quality control remains decisive: strength binds at $3.0\%$ air, and the batch-air penalty fails the targets; Section V-D distinguishes the codified and LOO-CV break-even limits. Reaching $\le2.0\%$ air requires PCE and 14--28\,d moist curing. The PCE route is uncalibrated on the physical cohort.

Allocation-invariant rankings and the uptake offsets support prioritizing clinker displacement over carbonation uptake. This conclusion remains subject to the exposure, strength, cover, allocation, and uncertainty conditions stated above; neither the local analytical branch nor the extrapolated mixes establish global or experimentally validated optimality.

\section*{Disclosure of AI Assistance}
ChatGPT (OpenAI) assisted with drafting parts of this research manuscript and refining its wording, clarity, and readability. The author retains responsibility for the manuscript's content and conclusions.

\bibliographystyle{IEEEtran}
\bibliography{references}

\appendices
\section{Standards and Test Details}
\label{sec:appendix_standards}
\textbf{Codes, Grades \& Protocols:} Exposure follows EN~206 \cite{en206}; the grade mapping, IS~10262 margins, and IS~456 limits are specified in Sections III-B and IV. PPC fly ash follows IS~1489 \cite{is1489}. Phenolphthalein ($1\%$) testing followed RILEM CPC-18 / EN~12390-10 \cite{hule2024}; strength and air-content testing followed EN~12390-3 \cite{en12390} and EN~12350-7 \cite{en12350}, respectively. Storage tests gave $p\ge0.55$; the durations are reported in Section II-A. Module B1 follows EN~16757 (Annex~BB) \cite{en16757} and Lagerblad \cite{lagerblad2005}. Achieving $\mathrm{Air}\le2.0\%$ at $w/b=0.45$ requires PCE, adding $+4.2\text{--}5.3\ \mathrm{kg\ CO_2e/m^3}$ \cite{efca2021}; the qualified carbon comparison is given in Table~\ref{tab:sensitivity_matrix}. No dataset mix used PCE, so this compaction route remains uncalibrated. Data: Mendeley Data \cite{hule2024mendeley} (\cite{hule2024}). Scripts: \url{https://github.com/abhitejj/low-carbon-concrete-optimization}. The study was conducted in \textit{MESS 2026}, under Prof.~R.~N.~Singh and TA Gayathri Vangala.

\clearpage
\onecolumn
\section{Detailed Mixtures, Screening Counts, and Analytical Derivations}
\label{sec:appendix_candidates}
\noindent Table~\ref{tab:candidates_details} retains all 18 scenarios from Table~\ref{tab:candidates_matrix}. Targets for C30/37, C35/45, and C40/50 are 30.0, 35.0, and 40.0 MPa. Continuous solutions use $w/b=0.45$ and a +8.25 MPa strength margin. Constituent percentages are by binder mass; a dash denotes zero replacement. Carbon is net GWP in kg CO$_2$e/m$^3$.

\begin{table}[!h]
\caption{Full Continuous Compositions and Discrete Feasibility Counts}
\label{tab:candidates_details}
\centering
\footnotesize
\renewcommand{\arraystretch}{1.15}
\setlength{\tabcolsep}{2.5pt}
\begin{tabular*}{\textwidth}{@{\extracolsep{\fill}}lc ccc ccccc ccccc@{}}
\toprule
& & \multicolumn{3}{c}{\textbf{Feasible discrete mixes}} & \multicolumn{5}{c}{\textbf{Continuous: baseline kinetics}} & \multicolumn{5}{c}{\textbf{Continuous: +15\% kinetics}} \\
\cmidrule(lr){3-5}\cmidrule(lr){6-10}\cmidrule(lr){11-15}
\textbf{Grade} & \textbf{Cover} & \textbf{Nom.} & \textbf{Kin.} & \textbf{Joint} & \textbf{FA} & \textbf{SG} & \textbf{LS} & \textbf{PO} & \textbf{Carbon} & \textbf{FA} & \textbf{SG} & \textbf{LS} & \textbf{PO} & \textbf{Carbon} \\
& \textbf{(mm)} & & & & \multicolumn{4}{c}{\textbf{(\% binder)}} & & \multicolumn{4}{c}{\textbf{(\% binder)}} & \\
\midrule
\multicolumn{15}{l}{\textbf{A. Sheltered (XC3): 35 candidates}} \\
C30/37 & 25 & 24/35 & 19/35 & 13 & 10.7 & 39.3 & -- & -- & \textbf{135.7} & 0.4 & 49.6 & -- & -- & \textbf{137.5} \\
C30/37 & 30 & 30/35 & 26/35 & 18 & 24.0 & 26.0 & -- & -- & \textbf{133.4} & 13.8 & 36.2 & -- & -- & \textbf{135.2} \\
C30/37 & 35 & 32/35 & 30/35 & 19 & 35.0 & 15.0 & -- & -- & \textbf{131.6} & 25.1 & 24.9 & -- & -- & \textbf{133.3} \\
\addlinespace[2pt]
C35/45 & 25 & 20/35 & 16/35 & 9 & 10.7 & 39.3 & -- & -- & \textbf{135.7} & 0.4 & 49.6 & -- & -- & \textbf{137.5} \\
C35/45 & 30 & 23/35 & 21/35 & 11 & 24.0 & 26.0 & -- & -- & \textbf{133.4} & 13.8 & 36.2 & -- & -- & \textbf{135.2} \\
C35/45 & 35 & 25/35 & 23/35 & 12 & 25.5 & 24.5 & -- & -- & \textbf{133.2} & 25.1 & 24.9 & -- & -- & \textbf{133.3} \\
\addlinespace[2pt]
C40/50 & 25 & 15/35 & 12/35 & 5 & 10.1 & 39.9 & -- & -- & \textbf{135.8} & 0.4 & 49.6 & -- & -- & \textbf{137.5} \\
C40/50 & 30 & 17/35 & 16/35 & 5 & 10.1 & 39.9 & -- & -- & \textbf{135.8} & 10.1 & 39.9 & -- & -- & \textbf{135.8} \\
C40/50 & 35 & 18/35 & 17/35 & 5 & 10.1 & 39.9 & -- & -- & \textbf{135.8} & 10.1 & 39.9 & -- & -- & \textbf{135.8} \\
\midrule
\multicolumn{15}{l}{\textbf{B. Unsheltered (XC4): 26 candidates}} \\
C30/37 & 25 & 24/26 & 22/26 & 14 & 35.0 & 12.2 & -- & 2.8 & \textbf{142.2} & 28.6 & 21.4 & -- & -- & \textbf{143.7} \\
C30/37 & 30 & 25/26 & 24/26 & 14 & 35.0 & 9.7 & 5.3 & -- & \textbf{142.0} & 35.0 & 10.4 & 2.5 & 2.1 & \textbf{142.0} \\
C30/37 & 35 & 25/26 & 25/26 & 14 & 35.0 & 9.7 & 5.3 & -- & \textbf{142.0} & 35.0 & 9.7 & 5.3 & -- & \textbf{142.0} \\
\addlinespace[2pt]
C35/45 & 25 & 17/26 & 16/26 & 9 & 25.5 & 24.5 & -- & -- & \textbf{144.2} & 25.5 & 24.5 & -- & -- & \textbf{144.2} \\
C35/45 & 30 & 17/26 & 17/26 & 9 & 25.5 & 24.5 & -- & -- & \textbf{144.2} & 25.5 & 24.5 & -- & -- & \textbf{144.2} \\
C35/45 & 35 & 17/26 & 17/26 & 9 & 25.5 & 24.5 & -- & -- & \textbf{144.2} & 25.5 & 24.5 & -- & -- & \textbf{144.2} \\
\addlinespace[2pt]
C40/50 & 25 & 13/26 & 13/26 & 5 & 10.1 & 39.9 & -- & -- & \textbf{146.9} & 10.1 & 39.9 & -- & -- & \textbf{146.9} \\
C40/50 & 30 & 13/26 & 13/26 & 5 & 10.1 & 39.9 & -- & -- & \textbf{146.9} & 10.1 & 39.9 & -- & -- & \textbf{146.9} \\
C40/50 & 35 & 13/26 & 13/26 & 5 & 10.1 & 39.9 & -- & -- & \textbf{146.9} & 10.1 & 39.9 & -- & -- & \textbf{146.9} \\
\bottomrule
\end{tabular*}
\end{table}

{\footnotesize
\noindent\textbf{Common exposure conditions.} Exposure bounds follow Section IV. ``Nom.'' counts nominal-strength/nominal-kinetics survivors; ``Kin.'' adds +15\% kinetic uncertainty alone; ``Joint'' applies both strength and kinetic buffering. Counts follow the same grade/cover order as Table~\ref{tab:candidates_matrix}.

\medskip
\noindent\textbf{Properties of selected discrete mixes.} C39 has measured strength 40.4 MPa and sheltered $d_{50}=22.6$ mm; C24 has measured strength 47.7 MPa, sheltered $d_{50}=25.04$ mm, and unsheltered $d_{50}=17.5$ mm.

\medskip
\noindent\textbf{Screening and admixture qualifications.} Nominal and buffered screening follow Section III-B. C39 also fails kinetic uncertainty alone ($25.99>25\ \mathrm{mm}$). The original mixes contained no PCE; adding 1.0\% PCE within the optimizer shifts benchmark Net GWP to $137.6\text{--}138.7\ \mathrm{kg/m^3}$. The $14.6\ \mathrm{kg/m^3}$ saving applies only to hypothetical identical PCE addition to C24, which requires none.

}


\subsection{Derivations for Closed-Form Durability and Boundary Optimality}
\label{sec:appendix_propositions}
{\footnotesize
\noindent\textbf{Durability reduction.} Substitution into Eq.~\ref{eq:kin_sh} gives 
\[
\begin{aligned}
\ln k_{\mathrm{sh}}&=-1.2939+4.0681(0.45)+0.02525p_{\mathrm{FA}}+0.01161(50-p_{\mathrm{FA}}) \\
&=1.1172+0.01364p_{\mathrm{FA}}.
\end{aligned}
\]
Enforcing $\ln k_{\mathrm{sh}}\le\ln(d_{\mathrm{allow}}/\sqrt{50})$ proves Proposition~1. At 30\,mm, $\ln(30/\sqrt{50})=1.4452$ gives $24.04\%$ FA and $25.96\%$ SG. Durability limits are $10.7\%$ FA at 25\,mm, $13.8\%$ at $26.09\ \mathrm{mm}$ (+15\% kinetics at 30\,mm cover), and $35.3\%$ at 35\,mm, reduced to $35.0\%$ by the PPC cap; strength may impose a tighter limit.

\medskip
\noindent\textbf{Objective slope and endpoint argument.} Under economic allocation, the binder slope is $\frac{B}{100}(EF_{\mathrm{FA}}-EF_{\mathrm{SG}})=-0.126\ \mathrm{kg\ CO_2e/(m^3\cdot\%FA)}$. With $\rho_{\mathrm{FA}}=2300<\rho_{\mathrm{SG}}=2900\ \mathrm{kg/m^3}$, aggregate displacement gives an embodied slope of $-0.130$; the benchmark uptake slope is $+0.037$, giving $\mathrm d\mathcal J/\mathrm dp_{\mathrm{FA}}=-0.167$ in the same units. Derivatives at $p_{\mathrm{FA}}=0,10,24.04,40,50\%$ are $-0.174,-0.172,-0.167,-0.158,-0.149$, respectively. These samples support monotonicity but do not prove its sign between samples. Conditional on a negative derivative throughout the feasible interval, increasing FA lowers $\mathcal J$ until durability, strength, or the composition cap binds, proving Proposition~2 on this fixed branch. Clinker emissions ($0.840\ \mathrm{kg/kg}$) exceed FA and SG ($0.015$ and $0.060$); this economic-allocation ordering alone does not prove global optimality.

\medskip
\noindent\textbf{Cover dividend and compaction arithmetic.} Where durability controls, $\mathrm dp_{\mathrm{FA}}^*/\mathrm dd_{\mathrm{allow}}=1/(0.01364d_{\mathrm{allow}})$ and the physical carbon dividend is $-(\mathrm d\mathcal J/\mathrm dp_{\mathrm{FA}})/(0.01364d_{\mathrm{allow}})$. At 30\,mm, $0.1672/(0.01364\times30)=0.4087\ \mathrm{kg\ CO_2e/(m^3\cdot mm)}$. KKT multiplier values depend on constraint scaling. The durability-controlled 25--35\,mm branch gives $(135.7-131.6)/10=0.41$, with interval slopes $0.46$ (25--30\,mm) and $0.36$ (30--35\,mm). These branch values apply while strength remains slack, as for C30/37; they are not the complete codified C35/45 path, which reaches $133.2\ \mathrm{kg/m^3}$ at 35\,mm. Benchmark strength slacks are $43.73-43.25=0.48\ \mathrm{MPa}$ codified, $4.33\ \mathrm{MPa}$ LOO-CV, and $8.73\ \mathrm{MPa}$ nominal. Break-even air is $2.0\%+0.48/4.354\approx2.11\%$ codified and $2.99\%\approx3.0\%$ LOO-CV.
}

\end{document}
